% Modernised reading — fonds Alexandre Grothendieck, shelfmark 156-4, pages 1-88.
%
% This is an INTERPRETATION, not a transcription. Everything the critical
% apparatus carried moves into footnotes. In any dispute of substance the
% transcription governs: whoever cites Grothendieck must cite the files
% batch-01.fr.tex to batch-05.fr.tex of this folder.
%
% Pass: Opus 5.5 (claude-opus-5-5), 2026-10-10 — first pass, unchecked against
% the pages by a human. The folder was read whole: five batches, pages 1-88,
% all mathematics, dated in his hand 10 June 1986 (p. 1), 11 June (p. 18),
% 12 June (p. 31), 13 June (p. 52), 14 June (p. 80).
%
% What the folder is. Chapter IV (« GF IV ») of Vers une géométrie des formes,
% « Analysis situs (première mouture) »: an attempt to axiomatise a « contrée »,
% a combinatorial model of a tame stratified space, from figures, their strata
% (multistrates), compatibility, refinement ≪ and subdivision ≼, down to the
% lieux (points) and the ombres of figures on them. The folder restarts its own
% framework five times (pp. 1, 25, 31-33, 45, 52); page 52 calls its own start
% « quatrième mouture ». These internal restarts are called « départs » here,
% and « mouture » is kept for the chapter-level drafts (156-4 the first, 156-6,
% 156-7, 156-8 the second to fourth).
%
% Notation fixed for the whole document. 𝕏 the ambient space; a figure F, its
% closed strata X_i, open strata U_i = X_i ∖ ∂X_i; Σ_F the ordered set of
% strata. « Multistrate » (his word from p. 19 on) for « figure élémentaire »
% (pp. 3-18) and « multiplicité » (pp. 21-37), the set 𝓜 for his 𝔉_él. ≤
% sub-figure / sub-multistrate, ≪ refinement, ≼ subdivision, R compatibility.
% 𝓛 the lieux; the ombre of X is 𝓛_X and its open ombre 𝓛°_X throughout (his
% X̃, X̃° on pp. 35-44, |X| on p. 87); 𝓜_{≪F} the multistrates refining F (his
% F̃ on p. 80); in the fifth start F̄ = 𝓜_{≤F} the set of strata. Axioms are
% always cited with their page, since « Cont 3 », « C 4 », « C 8 » each name
% several different statements.
%
% Substantive departures from the pages, each footnoted where it occurs:
% — p. 3: a sub-figure is DEFINED as a down-closed set of strata; the NB does not
%   follow from F' ⊂ F (his own counter-example, p. 27).
% — p. 6: a) ⇒ b) is proved directly (∂X_i = ∂X'_j) without the induction.
% — p. 9: « antifiltre » read as an ideal of sets; b) is illegible.
% — p. 14: ≪ is implied by ≤, i.e. weaker as a relation; « plus forte/fine ».
% — p. 18: the chain formula for geometric dimension counts one too many.
% — pp. 21, 35: existence of minimal strata needs Fig_2 / Cont'_3; supplied.
% — p. 26: the open strata cover |F| only for well-founded figures; example.
% — p. 28: the four conditions are not all equivalent: b') ⇔ d), and a) ⇔ b)
%   strictly weaker; example.
% — p. 33: the content of Cont 1 is the transitivity of ≤.
% — p. 41: F' ∪ F_X for his F ∪ F_X; c) ⇒ a) needs Cont 6 and Cont 7.
% — p. 47: C 4 as written (G' ∪ (F ∖ G)) is not closed; the margin's
%   correction is stated. Cor. 2 (≪ an order) is not proved there.
% — p. 51: the last step of the proof is not written.
% — p. 54: « Sup ↦ ∪ » and Cor. 2 a) need the strata to be sup-PRIME; the
%   diamond lattice M_3 satisfies C 1-C 2 otherwise. Supplied for pp. 54-68.
% — p. 68: the cosupport formula uses disjointness, not R.
% — p. 69: « modérée » read as « two subdivisions have a common one ».
% — p. 73: transitivity of ≪ is proved with C7_0, not C 6.
% — p. 76: G' for his F'. p. 78: the sign in b) inverted.
% — p. 80: the statement of the 14 June corollary is reconstructed from its
%   paraphrase and its proof (as read it is trivial).
% — p. 87: « contrées sans bords » is undefined; the converse fails for a circle.
%
% Modern names given and footnoted as ours: espace stratifié and condition de
% frontière (Whitney 1965, Thom 1969, Mather 1970), topologie d'Alexandrov
% (1937), topologie modérée / structures o-minimales (van den Dries; Pillay-
% Steinhorn 1986), idéal d'ensembles / bornologie, hauteur d'un ensemble
% ordonné, sous-polyèdre comme sous-complexe d'une subdivision, élément
% sup-irréductible et sup-premier, représentation de Birkhoff (1937) and Raney
% (1952), connexion de Galois / polarité (Birkhoff 1940), algèbre des ouverts
% réguliers d'un ensemble ordonné (forcing), espace extrêmement discontinu
% (Stone 1937, Gleason 1958).
%
% What the folder announces and never establishes: the converse of p. 8; that
% subdivisions are the maximal refinements in the topological model (p. 16
% stalls); the map F' → F of p. 42 (obtained only on p. 80, in the fifth
% start); that ≪ is an order on p. 47; the end of the proof of p. 51; that the
% supports form a Boolean algebra (p. 68, « je présume »); the Inf property
% of the induced refinement in the fifth start (Lemme 2, p. 78, abandoned);
% (*) of p. 83 for figures that are not multistrates; (ii)-(iii) of p. 85;
% the converse of p. 87; the Proposition of p. 88, on which the folder ends.
\documentclass[11pt,a4paper]{article}
\input{../preamble/grothendieck}

\folder{156-4}
\pages{1}{88}
\dating{1986}
\watermark{Édition de démonstration\\interprétation personnelle de l'œuvre}
\foldertitle{[Chapitre] IV. Analysis situs (première mouture) : notes manuscrites (10/06/1986) — lecture modernisée du dossier entier}

\begin{document}

\section*{Résumé}

\begin{resume}
Une carte de géographie découpe un pays en régions ; chaque région a des
frontières, les frontières se rencontrent en des points, et l'on peut toujours
redessiner la carte plus finement, en coupant une région en cantons, sans rien
changer au territoire. Ces quatre-vingt-huit feuillets, écrits du 10 au 14 juin
1986, cherchent à dire ce qu'est une telle carte quand on oublie le papier : à
décrire une « forme » géométrique uniquement par la façon dont ses morceaux
s'emboîtent, se touchent et se laissent recouper. Le mot qu'il choisit pour
l'objet ainsi décrit est justement celui d'un pays : une \emph{contrée}.

Le point de départ est un espace découpé en \emph{strates} : des morceaux
fermés dont l'intérieur est une variété connexe --- un sommet, une arête
ouverte, une face ouverte --- et dont le bord est fait de strates plus petites.
Un tel découpage est une \emph{figure} ; un morceau pris avec son bord est une
\emph{multistrate}. Deux opérations rendent une figure plus fine, et toute la
difficulté du dossier vient de ce qu'elles sont de nature différente. On peut
en prendre une partie (une \emph{sous-figure}, comme on ne garde que le contour
d'un triangle), ou la redécouper sans changer ce qu'elle recouvre (une
\emph{subdivision}, comme on coupe une arête par son milieu). Combiner les deux
donne le \emph{raffinement} : une sous-figure d'une subdivision --- c'est ainsi
qu'un petit segment tracé à l'intérieur d'un triangle « raffine » ce triangle.
Les plus fins de tous les objets sont les \emph{lieux}, qui ne se laissent plus
raffiner : dans l'exemple d'un espace, ce sont ses points. Toute figure jette
alors sur l'ensemble des lieux une \emph{ombre}, l'ensemble des points qu'elle
recouvre, et le dossier finit sur l'énoncé que cette ombre, avec la trace de
chaque strate, a précisément la forme d'une figure.

Il faut prévenir que le chemin n'est pas droit. En cinq jours, le cadre est
repris cinq fois --- « On reprend tout ! », « Donc, nouveau départ ! »,
« quatrième mouture ! » --- parce qu'à chaque fois un exemple (celui des
pseudo-disques, deux fois) montre que la notion prise pour primitive n'est pas
la bonne. Ce qui reste à la fin est une structure à deux ordres sur l'ensemble
des figures, l'un pour « être une partie de », l'autre pour « être un
redécoupage de », avec une liste d'axiomes qui disent comment ils se
répondent. Les marges notent « à vérifier », « il semblerait qu'il manque
quelque chose », « ouf ! », « J'abandonne ici ! ». Une phrase du 12 juin
dit la méthode mieux qu'aucun commentaire : plutôt que de dégager d'abord
l'axiomatique, dégager une terminologie géométrique, noter les axiomes au
passage, et n'axiomatiser qu'ensuite. Le dossier est cette terminologie en
train de se faire.

Ce que l'on y rencontre porte aujourd'hui d'autres noms : espaces stratifiés et
leur condition de frontière, ensembles ordonnés de strates, représentation d'un
treillis distributif par ses éléments irréductibles, connexions de Galois et
algèbres de Boole d'ouverts réguliers. Le chapitre suivant, « Algèbre des
figures » (dossier 156-5), commencé le 14 juin, reprend l'étude des familles de
parties sur lesquelles celui-ci s'arrête ; trois autres moutures de la même
Analysis situs suivent du 18 juin au 4 juillet (dossiers 156-6 à 156-8).

\keywords{stratified space, frontier condition, poset of strata, Alexandrov topology, subdivision, subcomplex of a subdivision, combinatorial dimension, join-irreducible element, Birkhoff representation theorem, compatibility relation, Galois connection, regular open algebra}
\end{resume}

\section*{Le fil du dossier, et les conventions}

Le dossier part d'un modèle --- un espace topologique stratifié --- et s'en
détache par étapes pour axiomatiser la structure qu'il y lit. Il le fait en
cinq départs successifs, qu'on garde distincts : chacun corrige le précédent,
et c'est ce mouvement qui est le contenu du dossier.

\begin{enumerate}
\item[I.] \emph{Premier départ : du modèle topologique à la contrée} (pages 1
à 24, 10 et 11 juin). Figures et strates d'un espace (pages 1 à 4) ;
compatibilité et configurations (pages 5 à 10) ; dimension combinatoire (pages
10 à 12) ; raffinement et subdivision (pages 13 à 16) ; la contrée abstraite et
ses lieux (pages 17 à 19) ; les ombres sur les lieux (pages 20 à 24).
\item[II.] \emph{Deuxième départ : les figures comme familles de parties des
lieux} (pages 25 à 31), ouvert par « On reprend tout ! ».
\item[III.] \emph{Troisième départ : multistrates, raffinement et compatibilité
primitifs} (pages 31 à 44, 12 juin), qui se clôt sur l'objection des
pseudo-disques.
\item[IV.] \emph{Quatrième départ : la subdivision primitive} (pages 45 à 51),
ouvert par « Donc, nouveau départ ! ».
\item[V.] \emph{Cinquième départ : l'ensemble ordonné des figures et ses deux
ordres} (pages 52 à 88, 13 et 14 juin), qu'il intitule « Contrée (quatrième
mouture !) »\footnote{Le mot « quatrième » est lu avec doute. Sa numérotation
ne compte vraisemblablement pas l'un des départs qu'on distingue ici ; on ne
cherche pas à deviner lequel.} : l'ordre $\leq$ seul (pages 52 à 68), puis la
subdivision $\preccurlyeq$ et le raffinement (pages 69 à 85), enfin les lieux
et les ombres (pages 86 à 88).
\end{enumerate}

\emph{Deux mots à ne pas confondre.} Le titre du dossier, « Analysis situs
(première mouture) », désigne le chapitre IV de \emph{Vers une géométrie des
formes} ; les chapitres VI, VII et VIII (dossiers 156-6, 156-7, 156-8) en sont
la deuxième, la troisième et la quatrième mouture. La « quatrième mouture » de
la page 52 est autre chose : un départ interne au présent dossier. On réserve
ici « mouture » aux chapitres et l'on dit « départ » pour les reprises
internes. De même, le \emph{lieu} de ce dossier --- une figure qu'on ne peut
plus raffiner, un point dans le modèle --- n'a rien à voir avec le lieu du
chapitre II (dossier 156-2), partie fermée qui coupe une faille en deux.

\emph{Conventions.} $\mathbb{X}$ est l'espace ambiant. Une figure $F$ a des
strates fermées $X_{i}$ et des strates ouvertes $U_{i} = X_{i} \smallsetminus
\partial X_{i}$ ; $\Sigma_{F}$ est l'ensemble ordonné de ses strates. On dit
\emph{multistrate} pour ce que les pages 3 à 18 appellent « figure
élémentaire » et les pages 21 à 37 « multiplicité » --- le mot est celui que la
page 19 adopte --- et l'on note $\mathcal{M}$ leur ensemble (son
$\mathfrak{F}_{\text{él}}$). Les relations sont notées : $\leq$ (sous-figure,
sous-multistrate), $\ll$ (raffinement), $\preccurlyeq$ (subdivision), $R$
(compatibilité). $\mathcal{L}$ est l'ensemble des lieux ; l'\emph{ombre} d'une
multistrate $X$, ensemble des lieux qui la raffinent, est notée
$\mathcal{L}_{X}$, et son \emph{ombre ouverte} $\mathcal{L}^{\circ}_{X}$, dans
tous les départs\footnote{Il les note $\mathcal{L}_{X}$ aux pages 20 à 24,
$\widetilde{X}$ et $\widetilde{X}^{\circ}$ aux pages 35 à 44, $|X|$ à la page
87.}. Les axiomes sont toujours cités avec leur page : « Cont 3 », « C 4 » ou
« C 8 » désignent chacun plusieurs énoncés différents selon le départ.

\emph{Ce que le dossier annonce et n'établit pas} : la réciproque de la page
8 ; que, dans le modèle topologique, les subdivisions soient exactement les
raffinements maximaux (la tentative de la page 16 s'arrête sur un problème de
topologie) ; l'application $F' \to F$ de la page 42, qui n'est obtenue qu'à la
page 80 ; que le raffinement soit une relation d'ordre au quatrième départ
(page 47) ; la dernière étape de la démonstration de la page 51 ; que les
supports forment une algèbre de Boole (page 68) ; la propriété d'Inf du
raffinement induit au cinquième départ (abandonnée page 78) ; les énoncés
(ii) et (iii) de la page 85 ; et la Proposition de la page 88, sur laquelle le
dossier s'arrête.

\pagerange{1}{4}
\section*{I. Le modèle topologique : figures et strates (pages 1 à 4)}

La page 1, datée du 10 juin 1986 et marquée « GF IV », annonce un modèle
« ensembliste topologique » de l'axiomatique d'un « modèle de formes
topologiques », dans le cas d'une contrée. L'espace ambiant $\mathbb{X}$ est
séparé, ou modéré pour une « théorie modérée » donnée, et dans ce cas « partie
fermée » s'entend « partie fermée modérée »\footnote{La « topologie modérée »
est celle que réclamait l'\emph{Esquisse d'un programme} (1984) ; voir la
lecture du chapitre II (dossier 156-2), pages 1 et 2. Elle a trouvé une forme
précise dans les structures o-minimales, développées dans les mêmes années
(L.~van den Dries ; A.~Pillay et C.~Steinhorn, 1986), sans lien avec ces
feuillets. Le rapprochement est le nôtre.}.

\textbf{Définition} (pages 1 et 2). Une \emph{figure} $F$ de $\mathbb{X}$ est
une partie fermée $X$, son \emph{support} $|F|$, munie d'une famille de parties
$X_{i}$, les \emph{strates fermées}, de réunion $X$, telles que :
\begin{enumerate}
\item[a)] les $X_{i}$ sont fermées, non vides et localement connexes ;
\item[b)] la famille des $X_{i}$ est localement finie ;
\item[c)] si $X_{i} \subsetneq X_{j}$, alors $X_{i}$ est rare dans $X_{j}$ ;
\item[d)] posant $\partial X_{i} = \bigcup_{X_{j} \subsetneq X_{i}} X_{j}$,
fermé par a), b), c), l'\emph{intérieur} $\mathrm{int}(X_{i}) = X_{i}
\smallsetminus \partial X_{i}$ est une variété connexe ;
\item[f)] $X_{i} \cap X_{j}$ est réunion de strates.
\end{enumerate}

On ordonne les strates par inclusion\footnote{Une condition e), écrite en
oblique dans la marge, en partie biffée, porte sur des dimensions ; elle n'est
lue que par fragments. La lettre f) suit donc d) sur la page.}. Comme
$\partial X_{i}$ est une réunion localement finie de fermés rares, elle est
rare, et la strate ouverte $U_{i} = \mathrm{int}(X_{i})$ est dense dans
$X_{i}$ : chaque strate fermée est l'adhérence de sa strate ouverte. Dans le cas
modéré, « variété » s'entend au sens modéré\footnote{Dans le cas purement
topologique, rien n'exige qu'une strate soit plongée de façon modérée : la
sphère cornue d'Alexander, avec la boule qu'elle borde, satisferait les
conditions. C'est une raison de plus pour l'hypothèse de modération de la page
1.}. Les strates ouvertes forment une partition de $|F|$ ; en termes actuels,
une figure est une \emph{stratification} localement finie d'un fermé de
$\mathbb{X}$ dont chaque strate fermée est réunion de strates --- ce qu'on
appelle la \emph{condition de frontière}\footnote{Les stratifications de
H.~Whitney (1965), R.~Thom (1969) et J.~Mather (1970) ajoutent des conditions
de régularité qui n'ont pas de sens ici, où rien n'est lisse. La condition
qu'une adhérence de strate soit réunion de strates est la « frontier
condition » de cette littérature. Les feuillets ne citent personne.}.

Une figure se reconstitue à partir de l'ensemble de ses strates : $F$ est un
élément de $\mathfrak{P}(\mathfrak{P}_{f}(\mathbb{X}))$, où
$\mathfrak{P}_{f}$ désigne les parties fermées, et l'ensemble $\mathfrak{F}$
des figures une partie de $\mathfrak{P}(\mathfrak{P}_{f}(\mathbb{X}))$. À
chaque figure est attaché l'ensemble ordonné $\Sigma_{F}$ de ses
strates\footnote{La page écrit « décomposition de $Y$ », « $\{X_{i}\} \in
\mathfrak{P}(Y)$ » pour la partie qu'elle vient d'appeler $X$.}.

\subsection*{Sous-figures et multistrates}

\begin{quote}
\textbf{Définition} (page 3). Une \emph{sous-figure} $F'$ de $F$ est une
partie de l'ensemble des strates de $F$ telle que, si $X_{i} \in F'$ et si
$X_{j}$ est une strate de $F$ contenue dans $X_{i}$, alors $X_{j} \in F'$. On
écrit $F' \leq F$.
\end{quote}

La page définit la sous-figure par la seule inclusion $F' \subset F$ et donne
la condition de clôture comme une conséquence ; elle n'en est pas
une\footnote{La page 27 donne elle-même le contre-exemple : un cercle $X_{1}$
muni d'un point $X_{0}$ forme une figure $F = \{X_{0}, X_{1}\}$, et $\{X_{1}\}$
est une figure contenue dans $F$, mais ce n'en est pas une sous-figure ---
l'intérieur de $X_{1}$ n'est pas le même dans les deux. On fait donc de la
clôture une partie de la définition.}. Ainsi définies, les sous-figures de $F$
sont en bijection avec les parties de $\Sigma_{F}$ fermées vers le bas ---
c'est-à-dire, comme le dit la page, avec les fermés de $\Sigma_{F}$ pour la
topologie dont les fermés sont les parties décroissantes\footnote{C'est la
topologie d'Alexandrov (P.~Alexandroff, 1937) de l'ensemble ordonné
$\Sigma_{F}$, à l'ordre près : le nom est le nôtre.} --- et $|F'| \subset
|F|$. La relation $\leq$ est un ordre sur $\mathfrak{F}$.

Une \emph{multistrate} est une figure qui a une strate maximale, c'est-à-dire
dont $\Sigma_{F}$ a un plus grand élément\footnote{La page dit « figure
élémentaire » ; une note de marge, avec la date « 11.6. », renvoie au mot
« multistrate » que la page 19 lui substituera. $\mathcal{M}$ rend son
$\mathfrak{F}_{\text{él}}$, cerclé sur la page.}. Les sous-multistrates d'une
figure $F$ correspondent à ses strates : à $X_{i}$ répond la multistrate
$F_{X_{i}}$ formée des strates de $F$ contenues dans $X_{i}$. Soit
$\Sigma'_{F} \subset \mathcal{M}$ l'ensemble de ces sous-multistrates ; on a
$\Sigma'_{F} \simeq \Sigma_{F}$, $F$ se reconstitue à partir de
$\Sigma'_{F}$, et l'inclusion des figures devient l'inclusion de ces parties de
$\mathcal{M}$. Ordonné par $\leq$, $\mathcal{M}$ est un ensemble ordonné, et
chaque $\Sigma'_{F}$ en est une partie fermée vers le bas.

\pagerange{5}{10}
\section*{II. Compatibilité et configurations (pages 5 à 10)}

Une partie de $|F|$ est \emph{$F$-saturée} si elle est réunion de strates
ouvertes de $F$ --- de façon équivalente, pour une partie fermée, réunion de
strates fermées\footnote{Le haut de la page 5, une première tentative de
caractériser les parties de $\mathcal{M}$ de la forme $\Sigma'_{F}$, est barré ;
une longue note de marge n'est lue que par fragments.}.

\textbf{Proposition} (pages 5 et 6). \emph{Pour deux figures $F$ et $F'$, les
conditions suivantes sont équivalentes :}
\begin{enumerate}
\item[a)] \emph{pour toutes strates $X \in \Sigma_{F}$, $X' \in \Sigma_{F'}$,
$X \cap X'$ est $F_{X}$-saturée et $F'_{X'}$-saturée ;}
\item[b)] \emph{deux strates ouvertes $U_{i}$ de $F$ et $U'_{j}$ de $F'$ sont
disjointes ou égales.}
\end{enumerate}
\emph{On dit alors que $F$ et $F'$ sont \textbf{compatibles}.}

b) $\Rightarrow$ a) : $X \cap X'$ est la réunion des $U_{k} \cap U'_{l}$, pour
$U_{k} \subset X$, $U'_{l} \subset X'$, chacune vide ou égale à $U_{k} =
U'_{l}$ ; c'est donc une réunion de strates ouvertes de l'une et de l'autre.
a) $\Rightarrow$ b) : soit $x \in U_{i} \cap U'_{j}$. La partie $X_{i} \cap
X'_{j}$ est fermée, $F_{X_{i}}$-saturée et contient $x$, donc $U_{i}$, donc son
adhérence $X_{i}$ ; ainsi $X_{i} \subset X'_{j}$, et symétriquement $X_{i} =
X'_{j}$. Reste à voir que $\partial X_{i} = \partial X'_{j}$ : si $Y$ est une
strate de $F$ strictement contenue dans $X_{i}$, a) appliqué à $(Y, X'_{j})$
dit que $Y$ est un fermé $F'_{X'_{j}}$-saturé distinct de $X'_{j}$ ; il ne
rencontre donc pas $U'_{j}$, sans quoi il contiendrait son adhérence $X'_{j}$,
et $Y \subset \partial X'_{j}$. D'où $\partial X_{i} \subset \partial X'_{j}$,
et l'inclusion inverse par symétrie\footnote{La page annonce une récurrence sur
$\sup(\dim F, \dim F')$, se ramenant aux figures $F_{X_{i}}$, $F'_{X'_{j}}$, et
s'interrompt ; le raisonnement est écrit vite et très corrigé. L'argument
direct ci-dessus est le nôtre ; il n'utilise que la densité des strates
ouvertes.}.

Deux figures sont compatibles si et seulement si leurs sous-multistrates le
sont deux à deux, et deux sous-figures d'une même figure le sont toujours
(page 7). On obtient sur $\mathcal{M}$ une relation réflexive et symétrique
$R$, « être compatibles ». Une \emph{configuration} est un ensemble de figures
deux à deux compatibles.

\subsection*{Les figures comme configurations}

Pour une figure $F$, la partie $\Phi = \Sigma'_{F}$ de $\mathcal{M}$ est
fermée vers le bas, c'est une configuration, et elle est localement finie : la
famille des supports de ses éléments est localement finie dans $\mathbb{X}$.
La page 8 affirme la réciproque : toute partie $\Phi$ de $\mathcal{M}$ ayant
ces trois propriétés provient d'une figure et d'une seule\footnote{Énoncée sans
démonstration. On la vérifie avec la Proposition des pages 5 et 6 : les strates
des éléments de $\Phi$, deux à deux compatibles, ont des strates ouvertes
disjointes ou égales, ce qui donne les conditions d) et f).}. D'où la
structure que la page 9 dégage, et qu'on numérote comme elle :
\begin{enumerate}
\item[(1)] l'ensemble $\mathcal{M}$ des multistrates ;
\item[(2)] l'ordre $\leq$ sur $\mathcal{M}$ ;
\item[(3)] la relation réflexive et symétrique $R$ de compatibilité ;
\item[(4)] parmi les configurations, celles qui sont « localement finies »,
$\mathrm{Conf}_{\mathrm{l.f.}} \subset \mathfrak{P}(\mathcal{M})$.
\end{enumerate}
La famille (4) contient les singletons, est stable par passage à une partie et
par réunion finie de configurations ; la page l'appelle un
\emph{antifiltre}\footnote{Le mot désigne ce qu'on appelle un idéal
d'ensembles, le dual d'un filtre ; une famille qui contient de plus les
singletons est une \emph{bornologie} (le nom est postérieur, H.~Hogbe-Nlend,
1977). La propriété b) de la page 9, qui lie la locale finitude d'une
configuration à celle de la partie fermée qu'elle engendre, est illisible en
son milieu.}. Les figures sont alors les parties de $\mathcal{M}$ qui sont des
configurations, fermées vers le bas et localement finies ; les \emph{figures
finies}, les parties finies qui sont des configurations fermées.

\pagerange{10}{12}
\section*{III. Dimension combinatoire (pages 10 à 12)}

Deux propriétés de la relation entre (2) et (3) :

\begin{quote}
\textbf{Fig}$_{1}$ (page 10). \emph{Deux sous-multistrates d'une même
multistrate $X$ sont compatibles}, de sorte que $\mathcal{M}_{\leq X}$ est une
configuration, donc une figure.

\textbf{Fig}$_{2}$ (page 11). \emph{Pour tout $X \in \mathcal{M}$, les chaînes
de $\mathcal{M}_{\leq X}$ ont un cardinal borné.}
\end{quote}

La \emph{dimension combinatoire} de $X$ est la borne supérieure des cardinaux
de ces chaînes, diminuée de un : c'est la hauteur de l'ensemble ordonné
$\mathcal{M}_{\leq X}$\footnote{Le nom de « hauteur » est le nôtre. Elle
compte les sauts de strate en strate, non la dimension des variétés : un
triangle fermé, décomposé en sommets, arêtes et face, a une dimension
combinatoire égale à $2$, mais un disque fermé dont le bord est une seule
strate circulaire a une dimension combinatoire égale à $1$.}. Elle est nulle
si et seulement si $X$ est minimal dans $\mathcal{M}$, et la dimension
combinatoire d'une figure est la borne supérieure de celles de ses éléments.
Dans le modèle, les multistrates minimales sont les fermés de $\mathbb{X}$ qui
sont des variétés connexes (une seule strate), et les figures de dimension
combinatoire nulle sont les configurations non vides et localement finies de
multistrates minimales, c'est-à-dire les fermés de $\mathbb{X}$ qui sont des
variétés, décomposés en leurs composantes connexes\footnote{La page écrit
« minimaux de $\mathfrak{F}$ » pour $\mathfrak{F}_{\text{él}}$, et
« décomposées » est lu avec doute. Ces sous-variétés ne sont pas supposées de
dimension constante, et, dans le cas topologique, aucune condition de
plongement n'est demandée.}. La figure vide correspond à la sous-variété
vide.

\pagerange{13}{16}
\section*{IV. Raffinement et subdivision (pages 13 à 16)}

\begin{quote}
\textbf{Définition} (page 14). Une figure $F'$ \emph{raffine} $F$, et l'on
écrit $F' \ll F$, si toute strate ouverte de $F'$ est contenue dans une strate
ouverte de $F$, nécessairement unique.
\end{quote}

On a alors $|F'| \subset |F|$, sans égalité nécessaire. C'est une relation
d'ordre --- deux figures qui se raffinent mutuellement ont même support et des
partitions en strates ouvertes plus fines l'une que l'autre, donc égales ---
et elle est impliquée par $\leq$, puisque les strates ouvertes d'une
sous-figure sont des strates ouvertes de la figure\footnote{La page dit de
$\ll$ qu'elle est « plus fine », puis « plus forte », que l'inclusion : elle
la contient, c'est-à-dire qu'elle relie plus de couples.}. Elle se lit sur les
multistrates : $F' \ll F$ si et seulement si toute sous-multistrate de $F'$
raffine une sous-multistrate de $F$. C'est la donnée
\begin{enumerate}
\item[(5)] l'ordre $\ll$ sur $\mathcal{M}$, impliqué par $\leq$,
\end{enumerate}
qui s'ajoute aux données (1) à (4), avec la relation (page 15)
\[
X \leq X' \iff X \ll X' \ \text{et}\ (X, X') \in R ,
\]
vraie dans le modèle : si chaque strate ouverte de $X$ est dans une strate
ouverte de $X'$ et lui est égale ou disjointe, elle lui est égale. Cette
équivalence --- sous-figure égale raffinement plus compatibilité --- revient
dans chacun des départs suivants.

\begin{quote}
\textbf{Définition} (pages 13 et 15). Une \emph{subdivision} de $F$ est un
raffinement $F'$ de même support.
\end{quote}

La page 15 annonce que cela revient à demander que $F'$ soit maximal, pour
$\leq$, parmi les raffinements de $F$, et en fait la définition
retenue\footnote{La phrase de la page 15 écrit « elle raffine $F'$ » pour
$F$. Une première version d'un point (6), biffée, cherchait à définir la
subdivision entre figures de $\mathcal{M}$.}. Une subdivision d'une
multistrate n'est pas nécessairement une multistrate : une arête coupée en son
milieu en a deux maximales.

La page 16, intitulée « Justification (?) », tente de montrer qu'un raffinement
$F'$ de support plus petit n'est pas maximal : elle forme la figure $F''$ dont
les strates ouvertes sont celles de $F'$ et les composantes connexes
$V_{i\alpha}$ des $U_{i} \smallsetminus X'$, où $X' = |F'|$. Il faut alors que
l'adhérence de chaque $V_{i\alpha}$ soit réunion de strates, et la page
s'arrête sur ce qu'elle appelle « un pb, qui a l'air intéressant, de topologie
d'un espace stratifié dans une variété »\footnote{C'est bien un problème : dans
un espace topologique quelconque, le complémentaire d'un fermé stratifié dans
une variété peut avoir une infinité de composantes, ou des composantes dont
l'adhérence n'est pas réunion de strates. Dans une structure o-minimale, toute
stratification définissable se raffine en une stratification compatible avec
un fermé définissable donné et satisfaisant la condition de frontière
(L.~van den Dries, \emph{Tame topology and o-minimal structures}, 1998) ; c'est
ce que la page demande. La longue note de marge de la page 16 n'est lue que par
fragments.}. L'équivalence reste ainsi annoncée et non établie.

\pagerange{17}{19}
\section*{V. La contrée abstraite et ses lieux (pages 17 à 19)}

La page 17 isole la structure. Une \emph{contrée} est donnée par :
\begin{enumerate}
\item[1°)] un ensemble $\mathcal{M}$ de multistrates ;
\item[2°)] un ordre $\ll$ sur $\mathcal{M}$, le raffinement ;
\item[3°)] une relation réflexive et symétrique $R$, la compatibilité ;
\item[4°)] un idéal de parties de $\mathcal{M}$, dites localement finies, qui
pour une contrée finie sont simplement les configurations finies,
\end{enumerate}
avec des axiomes « à dégager »\footnote{Le numéro 4°) est dans une note de marge
et n'est pas sûr ; plusieurs insertions obliques au-dessus de 1°) et 2°) ne sont
lues qu'en partie.}. Le premier : \textbf{Contr}$_{1}$ (page 17), \emph{la
relation « $X' \ll X$ et $X'$ compatible à $X$ » est un ordre} ; on la note
$\leq$, et Fig$_{1}$ se lit alors : deux multistrates qui raffinent $X$ et lui
sont compatibles sont compatibles entre elles. Le second :
\textbf{Contr}$_{2}$ (page 18), \emph{pour tout $X$, les dimensions
combinatoires des $X' \ll X$ sont bornées}\footnote{Une première version de
Contr$_{2}$, en haut de la page 18, est barrée.} ; leur borne supérieure est la
\emph{dimension géométrique} de $X$, et celle d'une figure est la borne
supérieure de celles de ses éléments. C'est donc le plus grand $n$ pour lequel
il existe une chaîne
\[
X_{0} < X_{1} < \cdots < X_{n}, \qquad X_{n} \ll X \in F
\]
(page 18)\footnote{La page écrit la chaîne $(X_{0}, \ldots, X_{n})$ avec $X_{0}
< \cdots < X_{n-1} \ll X_{n} \in F$ et en prend la longueur $n$. Comptée
ainsi, un segment serait de dimension $2$ (une extrémité, un sous-segment, le
segment) ; on compte les maillons de la chaîne stricte, ce qui redonne la
définition de la ligne précédente.}. Pour un segment, elle vaut $1$ ; pour
le disque dont le bord est une seule strate, la dimension combinatoire $1$ est
corrigée en dimension géométrique $2$ par ses subdivisions en cellules.

Le 11 juin, après un passage très raturé où il envisage de partir plutôt de
parties « constructibles », il annonce qu'il va « pousser un peu dans la
direction $\ll$, compatibilité »\footnote{Le bas de la page 18 est un palimpseste
d'insertions obliques ; on n'en donne que le fil principal, lui-même
fragmentaire.}, et arrive à la notion de lieu.

\begin{quote}
\textbf{Définition} (page 19). Un \emph{lieu} est une multistrate minimale pour
$\ll$, c'est-à-dire qui n'a d'autre raffinement qu'elle-même. On note
$\mathcal{L} \subset \mathcal{M}$ l'ensemble des lieux.
\end{quote}

Un lieu est a fortiori minimal pour $\leq$. Dans le modèle, les lieux sont
exactement les points de $\mathbb{X}$ : un point de $|X|$, pris comme figure à
une strate, raffine $X$, de sorte qu'une multistrate qui a plus d'un point
n'est pas minimale\footnote{La page conclut sur « les sous-variétés
$0$-connexes qui n'ont pas de raffinement autre qu'elles-mêmes » ; la lecture
du chiffre est sûre, celle du sens ne l'est pas. Ce sont les points.}. Les
figures sont désormais « des ensembles de multistrates mutuellement
compatibles, fermés pour $\leq$, localement finis ».

\pagerange{20}{24}
\section*{VI. Les ombres sur les lieux (pages 20 à 24)}

Pour une multistrate $X$, son \emph{ombre} $\mathcal{L}_{X}$ est l'ensemble
des lieux qui la raffinent, et $\widehat{\mathcal{L}}_{X} = \{\mathcal{L}_{Y}
\mid Y \leq X\} \subset \mathfrak{P}(\mathcal{L})$ ; pour une figure $F$,
$\mathcal{L}_{F} = \bigcup_{X \in F} \mathcal{L}_{X}$ et
$\widehat{\mathcal{L}}_{F} = \{\mathcal{L}_{X} \mid X \in F\}$. Le programme,
posé à la page 20, est de reconstituer $\ll$ et la compatibilité à partir de
$X \mapsto \widehat{\mathcal{L}}_{X}$\footnote{La page écrit
$\mathfrak{P}(\mathfrak{P}(X))$ là où l'on attend
$\mathfrak{P}(\mathfrak{P}(\mathcal{L}))$, et un « Contr$_{3}$ » y est commencé
puis barré ; à partir de la page 21, l'abréviation se lit « Cont ».}. Dans le
modèle, $\mathcal{L}_{X}$ est le support de $X$.

L'application $X_{i} \mapsto \mathcal{L}_{X_{i}}$ est croissante sur
$\mathcal{M}_{\leq X}$, et l'on veut qu'elle le soit strictement. On pose
l'\emph{ombre ouverte}
\[
\mathcal{L}^{\circ}_{X} = \mathcal{L}_{X} \smallsetminus \bigcup_{Y < X}
\mathcal{L}_{Y},
\]
et l'on demande :

\begin{quote}
\textbf{Cont 3} (page 21). \emph{Pour tout $X \in \mathcal{M}$,
$\mathcal{L}^{\circ}_{X} \neq \emptyset$.}

\textbf{Cont 4} (page 22). \emph{Si $X \neq Y$ sont compatibles,
$\mathcal{L}^{\circ}_{X} \cap \mathcal{L}^{\circ}_{Y} = \emptyset$.}
\end{quote}

Pour $x \in \mathcal{L}_{X}$, un élément $X_{i} \leq X$ minimal parmi ceux dont
l'ombre contient $x$ vérifie $x \in \mathcal{L}^{\circ}_{X_{i}}$, de sorte que
$\mathcal{L}_{X}$ est la réunion des $\mathcal{L}^{\circ}_{X_{i}}$ pour $X_{i}
\leq X$\footnote{L'existence d'un tel élément minimal n'est pas justifiée sur la
page ; elle résulte de Fig$_{2}$ (page 11), qui borne les chaînes de
$\mathcal{M}_{\leq X}$. On l'ajoute.}. Cont 4 dit que cette réunion est
disjointe.

\textbf{Scholie} (page 23). \emph{Soit $F$ une figure.}
\begin{enumerate}
\item[a)] \emph{Les ombres ouvertes $\mathcal{L}^{\circ}_{X}$, $X \in F$, sont
non vides et forment une partition de $\mathcal{L}_{F}$ ; en particulier
$X \mapsto \mathcal{L}^{\circ}_{X}$ est injective sur $F$.}
\item[b)] \emph{$X \mapsto \mathcal{L}_{X}$ est injective sur $F$, et pour
$X, Y \in F$ les conditions (i) $X \leq Y$, (ii) $\mathcal{L}_{X} \subset
\mathcal{L}_{Y}$, (iii) $\mathcal{L}^{\circ}_{X} \subset \mathcal{L}_{Y}$,
(iv) $\mathcal{L}^{\circ}_{X} \cap \mathcal{L}_{Y} \neq \emptyset$ sont
équivalentes.}
\item[c)] \emph{Pour $X, Y \in F$, $\mathcal{L}_{X} \cap \mathcal{L}_{Y} =
\bigcup_{Z \in F,\ Z \leq X,\ Z \leq Y} \mathcal{L}_{Z}$.}
\end{enumerate}

Pour (iv) $\Rightarrow$ (i) : $\mathcal{L}_{Y}$ est réunion des
$\mathcal{L}^{\circ}_{Y_{i}}$, $Y_{i} \leq Y$, et $\mathcal{L}^{\circ}_{X}$ en
rencontre un ; par a), $X = Y_{i}$. Les autres implications sont immédiates, et
seul sert le fait que $X$ et $Y$ sont compatibles, comme le note la marge.

Viennent les réciproques, posées en axiome :

\begin{quote}
\textbf{Cont 5} (page 24). \emph{$X \ll Y$ si et seulement si toute strate
ouverte de $X$ --- toute $\mathcal{L}^{\circ}_{X'}$, $X' \leq X$ --- est
contenue dans une strate ouverte de $Y$ ; $X$ et $Y$ sont compatibles si et
seulement si leurs strates ouvertes sont deux à deux disjointes ou égales.}
\end{quote}

On en déduit $\widehat{\mathcal{L}}_{X} = \widehat{\mathcal{L}}_{Y} \iff X = Y$,
et $X \leq Y$ si et seulement si $\widehat{\mathcal{L}}_{X}$ est une partie
fermée de $\widehat{\mathcal{L}}_{Y}$\footnote{Les signes d'égalité du a) du
Corollaire sont surchargés.}. La contrée $(\mathcal{M}, \ll, R)$ se reconstitue
donc à partir de $\mathcal{L}$ et de l'ensemble $\widehat{\mathcal{M}}$ des
$\widehat{\mathcal{L}}_{X}$, partie de
$\mathfrak{P}(\mathfrak{P}(\mathfrak{P}(\mathcal{L})))$ : c'est la conclusion du
premier départ, et la critique qu'en fera le 12 juin (page 31) porte
précisément sur elle.

\pagerange{25}{31}
\section*{VII. Deuxième départ : les figures comme familles de parties (pages 25 à 31)}

« On reprend tout ! » La conclusion de la page 24 est prise pour définition.

\begin{quote}
\textbf{Définition} (page 25). Une \emph{contrée} $C$ est un ensemble $L$, ses
\emph{lieux}, muni d'un ensemble $\mathfrak{F} \subset
\mathfrak{P}(\mathfrak{P}(L))$ de familles de parties de $L$, les
\emph{figures}. Les éléments d'une figure $F$ sont ses \emph{strates fermées},
leur réunion est son \emph{support} $|F|$.
\end{quote}

\begin{quote}
\textbf{Cont 1} (page 25). \emph{La famille vide est une figure, et pour tout
$x \in L$, $\{\{x\}\}$ est une figure, la figure ponctuelle.}

\textbf{Cont 2} (page 26). \emph{Pour toute figure $F$ et tout $X \in F$,
posant $\partial X = \bigcup_{Y \in F,\ Y \subsetneq X} Y$ et $X^{\circ} = X
\smallsetminus \partial X$, on a $X^{\circ} \neq \emptyset$ ; et pour $X, Y \in
F$, $X \cap Y$ est la réunion des $Z \in F$ contenus dans $X \cap Y$.}
\end{quote}

Une famille qui satisfait Cont 2 est dite \emph{admissible}\footnote{La note de
marge de la page 26 le dit, en partie lisible. C'est la notion que le chapitre
V, « Algèbre des figures » (dossier 156-5), commencé le 14 juin, prend pour
objet, avec sa variante « préadmissible ».}. Les strates ouvertes de deux
éléments distincts d'une famille admissible sont disjointes : si $x \in
X^{\circ} \cap Y^{\circ}$, un $Z \subset X \cap Y$ de $F$ contient $x$, et ne
peut être strictement contenu ni dans $X$ ni dans $Y$, d'où $X = Z = Y$. La
page en déduit que $X \mapsto X^{\circ}$ met $F$ en bijection avec les classes
d'une partition de $|F|$. Il faut pour cela que tout point de $|F|$ appartienne
à un élément minimal de $F$ parmi ceux qui le contiennent, ce que Cont 2
n'assure pas\footnote{Sur $L = \mathbf{N} \cup \{\omega\}$, la famille des
$X_{n} = \{\omega\} \cup \{n, n+1, \ldots\}$ satisfait Cont 2, avec $X_{n}^{\circ}
= \{n\}$, mais $\omega$ n'est dans aucune strate ouverte. L'hypothèse
manquante est automatique pour une figure finie, et plus généralement dès que
les chaînes décroissantes de $F$ sont stationnaires ; le chapitre V en fait la
condition $\beta$ de sa Proposition 1. On l'ajoute.}.

\begin{quote}
\textbf{Cont 3} (page 27). \emph{Toute partie d'une figure fermée vers le bas
pour l'inclusion est une figure} --- une \emph{sous-figure}, notée $F' \leq F$.
\end{quote}

C'est ici que la page note qu'une figure contenue dans une autre n'en est pas
nécessairement une sous-figure (voir la note de la page 3).

\subsection*{Compatibilité}

Pour deux figures $F$, $G$, la page 28 énumère des conditions qu'elle dit
équivalentes :
\begin{enumerate}
\item[a)] $|F| \cap |G|$ est saturée dans $|F|$ et dans $|G|$, et les deux
partitions qu'elle en reçoit coïncident ;
\item[b)] deux strates ouvertes de $F$ et de $G$ sont disjointes ou égales ;
\item[b')] pour $X \in F$, $Y \in G$, $X^{\circ} \cap Y^{\circ} = \emptyset$ ou
$X = Y$ ;
\item[d)] $F \cup G$ est admissible, et $F$, $G$ y sont fermées\footnote{Une
condition c), qui demande que les $X \cap Y$ soient saturés pour les
multistrates $\widehat{X}$, $\widehat{Y}$, n'est lue qu'en partie, et d)
elle-même est en partie illisible ; on la lit dans le sens que la suite de la
page lui donne.}.
\end{enumerate}
Ce n'est vrai qu'en partie. On a a) $\Leftrightarrow$ b), et b')
$\Leftrightarrow$ d) pour des figures dont tout point est dans une strate
ouverte (la note précédente) ; b') entraîne b), mais non la
réciproque\footnote{Sur $L = \{a, b, c\}$, les figures $F = \{\{a\}, \{a, b\}\}$
et $G = \{\{c\}, \{b, c\}\}$ satisfont a) et b) --- leurs strates ouvertes sont
$\{a\}, \{b\}$ et $\{c\}, \{b\}$ --- mais $\{a,b\} \cap \{b,c\} = \{b\}$ n'est
réunion d'aucun élément, et $F \cup G$ n'est pas admissible. Dans le modèle
topologique une strate fermée est l'adhérence de sa strate ouverte, et b) et
b') coïncident ; ici il n'y a pas de topologie, et une strate ouverte ne
détermine pas sa strate fermée. Les démonstrations de b') $\Leftrightarrow$ d)
reprennent l'argument du paragraphe précédent.}. On appelle donc
\emph{compatibles} deux figures qui satisfont d) ; deux sous-figures d'une même
figure le sont, et leur réunion en est une sous-figure\footnote{La page écrit
« sous-figures de $F \cup C$ » pour $F \cup G$.}. Une \emph{configuration} est
un ensemble de figures deux à deux compatibles.

\begin{quote}
\textbf{Cont 4} (page 29). \emph{Si toute multistrate de $F$ est compatible à
toute multistrate de $G$, alors $F$ et $G$ sont compatibles.}
\end{quote}

Ici la multistrate $\widehat{X}$ définie par $X \in F$ est la sous-figure des
éléments de $F$ contenus dans $X$. Par récurrence, une famille finie de figures
dont les multistrates sont deux à deux compatibles a pour réunion une figure,
dont elles sont des sous-figures\footnote{L'énoncé du corollaire de la page 29
est très surchargé ; seules les formules sont sûres.}.

La contrée est \emph{ensembliste}, ou \emph{contrée de points}, si deux figures
compatibles « ensemblistement » --- au sens de d) --- le sont, ce qui rend Cont
4 tautologique. Pour les autres, la page 30 propose :

\begin{quote}
\textbf{Cont}$'_{4}$ (pages 30 et 31). \emph{Deux figures sont compatibles si et
seulement si deux strates ouvertes $U$ de l'une et $V$ de l'autre sont égales
ou $C$-disjointes,}
\end{quote}
où deux lieux $x \neq y$ sont \emph{disjoints} quand $\{\{x\}, \{y\}\}$ est une
figure --- une $0$-sphère de $C$ ---, et deux parties $A$, $B$ de $L$ sont
\emph{$C$-disjointes} quand tout point de l'une est disjoint de tout point de
l'autre\footnote{L'énoncé de Cont$'_{4}$ à la page 30 est très surchargé, son
numéro aussi (un « 4 » et un « 1 » l'un sur l'autre) ; la forme citée est celle
de la reformulation du haut de la page 31.}. Dans une contrée de points, deux
lieux distincts sont toujours disjoints ; ailleurs, deux points distincts
peuvent être « trop proches » pour former une figure à eux deux --- c'est la
possibilité que ce départ ménage.

\pagerange{31}{34}
\section*{VIII. 12 juin : troisième départ (pages 31 à 40)}

« Depuis hier j'ai pas mal réfléchi à la construction des pseudo-disques
topologiques. » Il en ressort que si $(D, \partial D)$ est un pseudo-disque, les
lieux de $D \smallsetminus \partial D$ ne sont pas seulement des parties de
l'intérieur satisfaisant certaines conditions, mais doivent être accompagnés
d'une donnée supplémentaire liée au bord $\partial D$\footnote{La seconde moitié
de l'alinéa est écrite vite ; seuls les termes mathématiques sont sûrs. Le
pseudo-disque n'est pas défini dans le dossier, et l'on ne précise pas ici la
donnée supplémentaire, que la page ne dit pas.}. Une contrée ne peut donc pas
être une donnée de lieux et de familles de parties de lieux, du moins pour une
contrée « topologique combinatoire », et il faut prendre plutôt
\[
C = (\mathcal{M}, \ll, R_{\mathcal{M}}, \mathfrak{F}),
\]
multistrates, raffinement, compatibilité, et un ensemble $\mathfrak{F}$ de
parties de $\mathcal{M}$, les figures --- deux multistrates étant compatibles
quand elles appartiennent à une même figure, dit la marge.

Suit la phrase de méthode : plutôt que de dégager l'axiomatique, dégager une
terminologie géométrique des contrées, noter au passage les axiomes, et
n'axiomatiser qu'ensuite\footnote{Deux lignes très rapides ; la lecture n'en
retient que l'intention, qui est sûre.}. Une \emph{sous-contrée} est une partie
$\mathcal{M}'$ de $\mathcal{M}$ fermée pour $\ll$, munie des relations
induites, et la page prévient que la stabilité des axiomes par passage à une
sous-contrée n'est pas évidente.

\subsection*{Les axiomes « particulièrement cruciaux »}

\begin{quote}
\textbf{Cont 1} (page 33). \emph{La relation $X \leq Y$, définie par « $X \ll
Y$ et $X$ compatible à $Y$ », est une relation d'ordre.}
\end{quote}

Elle est réflexive et antisymétrique d'elle-même ; ce que l'axiome demande, c'est
la transitivité, c'est-à-dire que $X \leq Y \leq Z$ entraîne la compatibilité
de $X$ et de $Z$. Puis, pour les figures :
\begin{enumerate}
\item[Fig 1] pour $X \in \mathcal{M}$, $\mathcal{M}_{\leq X}$ est une figure,
la \emph{figure élémentaire} de $X$ ;
\item[Fig 2] les éléments d'une figure sont deux à deux compatibles, et une
figure est fermée vers le bas pour $\leq$ ;
\item[Fig 3] une partie fermée d'une figure est une figure ; la partie vide en
est une ;
\item[Fig 4] si tout élément de $F$ est compatible à tout élément de $G$,
$F \cup G$ est une figure, la \emph{figure réunion}\footnote{Les quatre axiomes
sont d'abord numérotés « Cont », numéros biffés et remplacés par « Fig » ; la
page écrit, dans Fig 1, « des $Y \in X$ » pour $Y \leq X$.}.
\end{enumerate}
Une contrée est \emph{finitiste} si ses figures finies sont exactement les
parties finies de $\mathcal{M}$, fermées pour $\leq$ et formées de multistrates
deux à deux compatibles ; elle est alors connue par $\ll$ et $R_{\mathcal{M}}$
seules, et la page 34 pose la question des conditions sur ces deux relations
pour que ces configurations satisfassent les axiomes\footnote{Que de telles
parties soient des figures se déduit de Fig 1 et Fig 4 par récurrence, dès que
la compatibilité passe aux sous-multistrates ; c'est le cas quand elle est
définie, comme dans la marge de la page 32, par l'appartenance à une même
figure. Un corollaire de la page 34 qui l'affirmait pour des raffinements
$X' \ll X$, $Y' \ll Y$ est barré : il serait faux.}.

\pagerange{35}{37}
\subsection*{Dimension, lieux, ombres (pages 35 à 37)}

\begin{quote}
\textbf{Cont}$'_{3}$ (page 35). \emph{Pour tout $X$, $\mathcal{M}_{\leq X}$ est
fini.}

\textbf{Cont 3} (page 35). \emph{Pour tout $X$, la dimension géométrique
$\sup_{X' \ll X} \dim \mathrm{comb}\, \mathcal{M}_{\leq X'}$ est finie.}
\end{quote}

Les \emph{lieux} sont les éléments de $\mathcal{M}$ minimaux pour $\ll$ ;
l'ombre $\mathcal{L}_{X}$ et l'ombre ouverte $\mathcal{L}^{\circ}_{X}$ sont
définies comme à la page 21, et l'ombre d'une figure est $\mathcal{L}_{F} =
\bigcup_{X \in F} \mathcal{L}_{X}$. Les axiomes Cont 3 et Cont 4 des pages 21
et 22 reviennent sous le nom de \textbf{Cont 4} (page 35) : a) les ombres
ouvertes sont non vides ; b) celles de deux multistrates compatibles distinctes
sont disjointes.

\textbf{Scholie} (pages 36 et 37). \emph{Soit $F$ une figure.}
\begin{enumerate}
\item[a)] \emph{$X \mapsto \mathcal{L}_{X}$ est injective sur $F$, et son image
est une famille admissible de parties de $\mathcal{L}$ au sens de Cont 2 (page
26) ;}
\item[b)] \emph{les $\mathcal{L}^{\circ}_{X}$, $X \in F$, forment une partition
de $\mathcal{L}_{F}$ indexée par $F$ ;}
\item[c)] \emph{les conditions (i) à (iv) de la page 23 sont équivalentes ;}
\item[d)] \emph{$F' \mapsto \mathcal{L}_{F'}$ met les sous-figures de $F$ en
bijection avec les parties de $\mathcal{L}_{F}$ fermées pour la famille
admissible des $\mathcal{L}_{X}$.}
\end{enumerate}

C'est le Scholie de la page 23, avec cette précision que l'ombre ouverte de $X$
est la strate ouverte de $\mathcal{L}_{X}$ dans la famille admissible --- d'où
les noms d'\emph{ombre ouverte} et de \emph{bord d'ombre} que la page leur
donne. Pour que tout lieu de $\mathcal{L}_{X}$ soit dans une ombre ouverte, il
faut un élément minimal, que Cont$'_{3}$ fournit\footnote{Comme à la page 21 ;
on l'ajoute. La formule du d) est surchargée et sa lecture douteuse.}. La
Proposition de la page 37 vérifie qu'une sous-contrée $\mathcal{M}'$, munie des
figures de $\mathcal{M}$ contenues dans $\mathcal{M}'$, satisfait encore Cont 1
à Cont 4, Fig 1 à Fig 4 et Cont$'_{3}$, avec les mêmes dimensions et les mêmes
ombres, et qu'une sous-contrée d'une contrée finitiste l'est aussi.

\pagerange{38}{40}
\subsection*{Disjonction, contrées induites, subdivision (pages 38 à 40)}

Deux multistrates $X$, $Y$ sont \emph{disjointes} si elles sont compatibles
et n'ont aucune sous-multistrate commune ; deux figures, si elles sont
compatibles et d'intersection vide. Une figure $F$ \emph{raffine} une figure
$G$ si tout élément de $F$ raffine un élément de $G$\footnote{La page note
cette relation entre figures par un triple chevron, $F \lll G$ ; on garde $\ll$,
qui coïncide avec elle sur les figures élémentaires.}.

\begin{quote}
\textbf{Proposition} (page 38). \emph{$F$ est une sous-figure de $G$ si et
seulement si $F$ raffine $G$ et lui est compatible.}
\end{quote}

En effet, pour $X \in F$ il existe $Y \in G$ avec $X \ll Y$, et $X$, $Y$ sont
compatibles, donc $X \leq Y$ et $X \in G$ par Fig 2.

\begin{quote}
\textbf{Cont 5} (page 39). \emph{Si $X$ et $Y$ sont disjointes, tout
raffinement de $X$ est disjoint de tout raffinement de $Y$.}

\textbf{Cont 6} (page 39). \emph{Si $\mathcal{L}_{X} \cap \mathcal{L}_{Y} =
\emptyset$ et si tout lieu de $\mathcal{L}_{X}$ est compatible à tout lieu de
$\mathcal{L}_{Y}$, alors $X$ et $Y$ sont disjointes.}
\end{quote}

Cont 5 dit en particulier que deux multistrates disjointes n'ont aucun
raffinement commun, et il s'étend aux figures : deux raffinements de figures
disjointes sont disjoints (Corollaire, page 39). Cont 6 dit que la disjonction
se lit sur les lieux. Les deux
passent aux sous-contrées. Pour une figure $F$, la \emph{contrée induite}
$C_{F}$ est formée des multistrates qui raffinent $F$, la \emph{contrée
résiduelle} $C_{\complement F}$ de celles qui sont disjointes de $F$, fermée
pour $\ll$ grâce à Cont 5\footnote{La page définit encore une contrée
résiduelle « de $F$ dans $G$ », $C_{F} \cap C_{\complement G}$, dont les lettres
sont surchargées ; l'exposant de $\mathcal{M}^{\sharp}_{\ll}$ est douteux.}.

\begin{quote}
\textbf{Définition} (page 40). Une \emph{subdivision} de $F$ est un raffinement
de $F$ maximal pour l'inclusion parmi les raffinements de $F$.
\end{quote}

C'est la définition de la page 15, transportée dans la contrée abstraite.

\pagerange{41}{44}
\section*{IX. Ce que le raffinement ne permet pas (pages 41 à 44)}

\textbf{Proposition} (page 41). \emph{Supposons Cont 6 et Cont 7 ci-dessous.
Pour un raffinement $F'$ de $F$, sont équivalents :}
\begin{enumerate}
\item[a)] \emph{$F'$ est une subdivision de $F$ ;}
\item[b)] \emph{aucune multistrate qui raffine $F$ n'est disjointe de $F'$ ;}
\item[c)] \emph{aucun lieu de $\mathcal{L}_{F}$ n'est disjoint de $F'$.}
\end{enumerate}

a) $\Rightarrow$ b) : si $X \ll F$ est disjointe de $F'$, la figure $F' \cup
\mathcal{M}_{\leq X}$ raffine $F$ et contient strictement $F'$\footnote{La page
écrit $F \cup F_{X}$.}. b) $\Rightarrow$ c) est immédiat, les lieux étant des
multistrates. c) $\Rightarrow$ a) : si $F' \subsetneq F'' \ll F$, choisissons
$X \in F'' \smallsetminus F'$ et $x \in \mathcal{L}^{\circ}_{X}$ ; alors $x \in
\mathcal{L}_{F}$, et il faut voir que $x$ est disjoint de $F'$. La page s'arrête
là-dessus --- « il faut un axiome pour l'assurer » --- et pose :

\begin{quote}
\textbf{Cont 7} (page 41). \emph{Si $X \neq Y$ sont compatibles, tout lieu de
$\mathcal{L}^{\circ}_{X}$ est disjoint de tout lieu de
$\mathcal{L}^{\circ}_{Y}$.}
\end{quote}

Avec lui, $x$ est disjoint de tous les lieux des ombres ouvertes des éléments
de $F'$, qui recouvrent $\mathcal{L}_{F'}$, et Cont 6 conclut que $x$ est
disjoint de $F'$. C'est aussi le Corollaire 1 de la page 41 : si $F'$ est une
sous-figure de $F$ et $X \in F \smallsetminus F'$, tout lieu de
$\mathcal{L}^{\circ}_{X}$ est disjoint de $F'$\footnote{La page invoque Cont 6 entre parenthèses et note « à
vérifier » en marge ; l'enchaînement est le nôtre. La page 52 résume Cont 7 en
une phrase : « deux lieux de strates ouvertes différentes sont disjoints ». Dans le Corollaire 1, « disjoint » est écrit au-dessus de « compatible » biffé, et lu avec doute.}.

\subsection*{Une application qui manque, et le raffinement induit}

La page 42 veut associer à chaque élément $X'$ d'un raffinement $F'$ de $F$
l'unique $X \in F$ dont l'ombre ouverte contient celle de $X'$ --- la version,
dans la contrée, de la définition topologique de la page 14 ---, prend pour
$X$ un élément minimal de $F$ raffiné par $X'$, et ne parvient pas à montrer
que $\mathcal{L}^{\circ}_{X'} \subset \mathcal{L}^{\circ}_{X}$ : « Il semblerait
qu'il manque quelque chose. » Ce qui manque est posé à la page 43 :

\begin{quote}
\textbf{Cont 8} (page 43). \emph{Soient $Y \ll X$ et $X' < X$. Les $Y' \leq Y$
qui raffinent $X'$ forment une sous-figure $R_{X'}$ de $\mathcal{M}_{\leq Y}$,
le \textbf{raffinement de $X'$ induit par $Y$}, et tout raffinement commun de
$Y$ et de $X'$ raffine $R_{X'}$ : $R_{X'}$ est la borne inférieure de
$\mathcal{M}_{\leq Y}$ et de $\mathcal{M}_{\leq X'}$ pour $\ll$.}
\end{quote}

En particulier $\mathcal{L}_{Y} \cap \mathcal{L}_{X'} = \bigcup_{Y' \in R_{X'}}
\mathcal{L}_{Y'}$, et, s'il n'existe pas de $X_{1}$ avec $Y \ll X_{1} < X$, on
a $\mathcal{L}^{\circ}_{Y} \subset \mathcal{L}^{\circ}_{X}$ : un lieu $x$ de
$\mathcal{L}^{\circ}_{Y}$ qui serait dans $\mathcal{L}_{X'}$, $X' < X$, serait
dans l'ombre d'un $Y' \in R_{X'}$, donc $Y' = Y$ par définition de l'ombre
ouverte, et $Y \ll X' < X$. D'où, en prenant pour chaque strate de $Y$ la plus
petite strate de $X$ qu'elle raffine :

\begin{quote}
\textbf{Corollaire} (page 44). \emph{Si $F'$ raffine $F$, toute ombre ouverte
d'un élément de $F'$ est contenue dans l'ombre ouverte d'un unique élément de
$F$.}
\end{quote}

\subsection*{L'objection des pseudo-disques}

« Je viens de m'apercevoir » que, dans le cas des pseudo-disques, deux lieux
disjoints $x$, $y$ de l'intérieur ne doivent pas être des \emph{nœuds} d'une
même subdivision --- les nœuds d'une figure étant ses éléments qui sont des
lieux, $F \cap \mathcal{L}$. La figure $\{x, y\}$, dont chaque élément raffine
le disque, ne doit donc pas compter comme un raffinement, ou en tout cas ne
définit pas de subdivision\footnote{La construction de la première phrase est
celle de la page, et sa lecture est en partie douteuse ; on en garde la
conclusion, qui est sûre.}. Il est alors inapproprié de dire que $F'$ raffine
$F$ quand chaque élément de $F'$ raffine un élément de $F$ : il vaut mieux
prendre pour primitive la \emph{subdivision} d'une figure et définir un
\emph{raffinement} comme une sous-figure d'une subdivision\footnote{C'est la
caractérisation classique des sous-polyèdres en topologie linéaire par morceaux
: une partie d'un polyèdre $|K|$ en est un sous-polyèdre si et seulement si
c'est un sous-complexe d'une subdivision de $K$ (voir C.~Rourke et
B.~Sanderson, 1972). Le rapprochement est le nôtre.}. Toute la suite du dossier
est bâtie sur ce renversement.

\pagerange{45}{51}
\section*{X. Quatrième départ : la subdivision primitive (pages 45 à 51)}

« Donc, nouveau départ ! » Une contrée est un triple $(\mathcal{M}, \mathfrak{F},
\Sigma)$ : un ensemble de multistrates, un ensemble $\mathfrak{F} \subset
\mathfrak{P}(\mathcal{M})$ de figures, et une relation $\Sigma$ sur
$\mathfrak{F}$, « $F'$ subdivise $F$ ».

\begin{quote}
\textbf{C 1} (page 45). \emph{Toute $X \in \mathcal{M}$ appartient à une plus
petite figure $F_{X}$, et $X \neq Y$ entraîne $F_{X} \neq F_{Y}$.} On pose
$X \leq Y$ si $F_{X} \subset F_{Y}$, c'est-à-dire $X \in F_{Y}$, et l'on veut
que ce soit un ordre.
\end{quote}

Toute figure est alors fermée vers le bas pour $\leq$. Deux figures sont
\emph{compatibles} si leur réunion est une figure, deux multistrates si leurs
figures élémentaires le sont.

\begin{quote}
\textbf{C 2} (page 46). \emph{$F$ et $G$ sont compatibles si et seulement si
tout élément de $F$ est compatible à tout élément de $G$.}

\textbf{C 3} (page 46). \emph{Toute partie d'une figure fermée pour $\leq$ est
une figure.}
\end{quote}

En particulier la figure vide existe dès que $\mathcal{M} \neq \emptyset$, et
on la suppose aussi quand $\mathcal{M} = \emptyset$ (la \emph{contrée vide}).

\begin{quote}
\textbf{Définition} (page 47). $G$ \emph{raffine} $F$, $G \ll F$, s'il existe une
subdivision $F'$ de $F$ dont $G$ est une sous-figure.
\end{quote}

Pour que ce soit un ordre, la page pose un axiome C 4 sous une forme qu'elle
corrige aussitôt en marge (« Corriger », puis « OK ») :

\begin{quote}
\textbf{C 4} (page 47, forme corrigée). \emph{Soient $G$ une sous-figure de $F$
et $G'$ une subdivision de $G$. Il existe une subdivision de $F$ qui induit
$G'$ sur $G$, et $G' \cap F \subset G$.}
\end{quote}

La forme écrite dans le corps de la page --- « $G' \cup (F \smallsetminus G)$
est une subdivision de $F$ » --- ne peut pas être gardée : $F \smallsetminus G$
n'est pas fermée en général\footnote{Si $F$ est un triangle et $G$ l'une de ses
arêtes, coupée en deux par $G'$, la multistrate « face » de $F \smallsetminus
G$ a pour bord l'arête entière, qui n'est plus dans $G' \cup (F \smallsetminus
G)$. Une subdivision de $F$ qui induit $G'$ doit remplacer aussi la face par
une face dont le bord est subdivisé ; c'est ce que dit la correction de la
marge.}. Les corollaires notés « OK » en marge sont : toute subdivision d'une
sous-figure est induite par une subdivision de $F$, et $\ll$ est une relation
d'ordre. Le second n'est pas démontré, et ne peut pas l'être avec ce qui
précède : l'antisymétrie de $\ll$ demandera un axiome (pages 73 à 83).

\begin{quote}
\textbf{Corollaire 3} (page 48). \emph{$G$ est une sous-figure de $F$ si et
seulement si $G$ est compatible à $F$ et la raffine.}
\end{quote}

Si $G \subset F'$, subdivision de $F$, et $H = F \cup G$ est une figure, la
forme corrigée de C 4, appliquée à la sous-figure $F$ de $H$ et à sa subdivision
$F'$, donne $G \subset F' \cap H \subset F$. Pour une multistrate, on retrouve
l'équivalence des pages 15, 17 et 33 : $X \leq Y$ si et seulement si $X \ll Y$
et $X$ compatible à $Y$.

Deux figures sont \emph{disjointes} si elles sont compatibles et d'intersection
vide, deux multistrates si leurs figures élémentaires le sont. Les axiomes
suivants règlent les rapports de la subdivision avec le reste :

\begin{quote}
\textbf{C 5} (page 49). \emph{Si $F'$ subdivise $F$, une figure est disjointe
de $F$ si et seulement si elle l'est de $F'$.}

\textbf{C 6} (page 49). \emph{Si $F'$ subdivise $F$ et $G \leq F$, il existe
une unique subdivision $G'$ de $G$ contenue dans $F'$, la \textbf{subdivision
induite} $F'|G$, formée des $X \in F'$ qui raffinent $G$.}

\textbf{C 7} (pages 49 et 50). \emph{Si $F = \bigcup F_{i}$ est réunion de
sous-figures et si les subdivisions $F'_{i}$ de $F_{i}$ induisent les mêmes
subdivisions sur les $F_{i} \cap F_{j}$, il existe une unique subdivision $F'$
de $F$ avec $F'|F_{i} = F'_{i}$ pour tout $i$.}
\end{quote}

Par C 5 et C 4, la réunion d'une subdivision de $F$ et d'une figure disjointe
$G$ subdivise $F \cup G$, et des raffinements de figures disjointes sont
disjoints\footnote{La page écrit $G' \ll F$ pour $G' \ll G$.} ; par C 7, deux
subdivisions de $F$ qui induisent la même subdivision sur chaque multistrate
sont égales.

\begin{quote}
\textbf{Proposition} (page 50). \emph{Soient $G \leq F$ et $F' \ll F$. Alors
$G' = \{X \in F' \mid X \ll G\}$ est une sous-figure de $F'$ qui raffine $G$,
et c'est la plus grande sous-figure de $F'$ qui raffine $G$.}

\textbf{Corollaire} (page 51). \emph{$G'$ est la borne inférieure de $F'$ et
de $G$ pour $\ll$ : toute figure qui raffine $F'$ et $G$ raffine $G'$.}
\end{quote}

Pour la Proposition, il y a quelque chose à prouver, puisque des éléments qui
raffinent $G$ ne forment pas pour autant un raffinement de $G$. Soit
$\overline{F}'$ une subdivision de $F$ contenant $F'$, et $\overline{G}' =
\overline{F}'|G$ (C 6) ; alors $G' = F' \cap \overline{G}'$ est une
sous-figure de $\overline{G}'$, donc raffine $G$\footnote{La page écrit « la plus
grande sous-figure de $F'$ qui raffine $G'$ » pour $G$.}. Le Corollaire est
démontré à la page 51 au moyen de quatre schémas de subdivisions successives
et de C 4 « revu et corrigé », et se termine par « qed, ouf ! (C'était
l'ancien Cont 8) » --- c'est en effet l'axiome Cont 8 de la page 43, devenu
théorème\footnote{La démonstration établit qu'une figure $\Phi$ qui raffine
$F'$ et $G$ est contenue dans la subdivision $\overline{G}'$ de $G$ ; le pas
qui en tire $\Phi \ll G'$ n'est pas écrit, et on ne le refait pas ici. Les
schémas de la page ne sont pas redessinés.}.

\pagerange{52}{55}
\section*{XI. 13 juin : cinquième départ, l'ensemble ordonné des figures (pages 52 à 60)}

Le 13 juin, il fait le compte des anciens axiomes qui restent à exprimer --- la
dimension finie (Cont 3), l'admissibilité (Cont 4), la disjonction lue sur les
lieux (Cont 6), Cont 7 --- puis : « avant de continuer, je vais encore une fois
changer d'optique ! »

\begin{quote}
\textbf{Définition} (page 52). Une \emph{contrée} est un ensemble
$\mathfrak{F}$ de \emph{figures} muni de deux relations d'ordre, $\leq$
(\emph{sous-figure}, notée aussi $\subset$) et $\preccurlyeq$
(\emph{subdivision}) :
\[
\mathcal{C} = (\mathfrak{F}, \leq, \preccurlyeq).
\]
\end{quote}

On commence par $\leq$ seul.

\begin{quote}
\textbf{C 1} (page 53). a) $\mathfrak{F} \neq \emptyset$ ; b) \emph{toute
famille majorée de figures a une borne supérieure.}
\end{quote}

La famille vide est majorée, d'où un plus petit élément, la \emph{figure vide}
$\emptyset_{\mathfrak{F}}$. Deux figures sont \emph{compatibles} si elles sont
majorées, et leur borne supérieure est leur \emph{réunion} $F \cup G$ ; une
\emph{configuration} est une famille majorée. Une figure est
\emph{irréductible} si $F = \mathrm{Sup}_{i} F_{i}$ entraîne $F = F_{i}$ pour un
$i$ ; la figure vide, borne supérieure de la famille vide, ne l'est pas. Les
figures irréductibles sont les \emph{multistrates}, leur ensemble est
$\mathcal{M}$, et les multistrates majorées par $F$ sont ses \emph{strates},
d'ensemble $\overline{F}$\footnote{En termes actuels : les éléments
complètement sup-irréductibles de l'ensemble ordonné $(\mathfrak{F}, \leq)$.
Le nom est le nôtre.}.

\begin{quote}
\textbf{C 2} (page 54). \emph{Toute figure est la borne supérieure de ses
strates.}
\end{quote}

L'application $F \mapsto \overline{F}$, de $\mathfrak{F}$ dans
$\mathfrak{P}(\mathcal{M})$, est alors injective et identifie $(\mathfrak{F},
\leq)$ à l'ensemble $\Phi$ des $\overline{F}$ ordonné par inclusion. La page
ajoute qu'elle transforme les bornes supérieures en réunions, et que $\Phi$ est
donc stable par réunion de familles majorées. Pour cela il faut une hypothèse
que C 1 et C 2 ne donnent pas : qu'une strate de $\mathrm{Sup}_{i} F_{i}$ soit
majorée par l'un des $F_{i}$, autrement dit que les multistrates soient
\emph{sup-premières}\footnote{Le treillis « diamant » $M_{3}$ --- un plus petit
élément, trois atomes $a$, $b$, $c$, un plus grand élément $1$ --- satisfait C 1
et C 2, ses irréductibles étant $a$, $b$, $c$ ; mais $\overline{a \vee b} =
\{a, b, c\} \neq \{a\} \cup \{b\}$. Avec l'hypothèse, $(\mathfrak{F}, \leq)$ est
de la forme décrite à la page 57 ; c'est l'analogue, pour ces ensembles
ordonnés, de la représentation d'un treillis distributif fini par l'ensemble
ordonné de ses sup-irréductibles (G.~Birkhoff, 1937 ; G.~Raney, 1952, pour le
cas infini). L'hypothèse vaut dans le modèle topologique, où sous-figures et
réunions sont ensemblistes. Les feuillets ne citent personne.}. On la suppose
dans tout ce qui suit, jusqu'à la page 68.

Réciproquement (page 55), si $\mathcal{M}$ est un ensemble et $\Phi \subset
\mathfrak{P}(\mathcal{M})$ une famille non vide, stable par réunion de familles
majorées, telle que tout $X$ appartienne à un plus petit $A_{X} \in \Phi$, alors
$(\Phi, \subset)$ satisfait C 1 et C 2, ses irréductibles sont les $A_{X}$, et
$X \mapsto A_{X}$ identifie $\mathcal{M}$ à l'ensemble de ses multistrates ; une
figure s'identifie alors à l'ensemble de ses strates. On ordonne
$\mathcal{M}$ par $\leq$, « jamais par $\preccurlyeq$ ! ».

\pagerange{56}{60}
\subsection*{Sous-figures, lieux, disjonction, compatibilité (pages 56 à 60)}

\begin{quote}
\textbf{Proposition} (page 56). a) \emph{Les sous-figures de $F$
s'identifient, par $F' \mapsto \overline{F}'$, aux parties de $\overline{F}$
fermées vers le bas.} b) \emph{Deux figures, et plus généralement une famille
non vide de figures, ont une borne inférieure, et $\overline{F \wedge G} =
\overline{F} \cap \overline{G}$.}
\end{quote}

Les \emph{lieux} sont les figures minimales parmi les figures non vides ; ils
sont irréductibles, et ce sont aussi les éléments minimaux de $\mathcal{M}$.
Deux figures sont \emph{disjointes} si elles sont compatibles et si $F \wedge G
= \emptyset_{\mathfrak{F}}$.

\begin{quote}
\textbf{Proposition} (pages 56 et 57). \emph{Si $F$ et $G$ sont disjointes,
$(F', G') \mapsto F' \cup G'$ est une bijection}
\[
\mathfrak{F}_{\leq F} \times \mathfrak{F}_{\leq G} \xrightarrow{\ \sim\ }
\mathfrak{F}_{\leq F \cup G},
\]
\emph{d'inverse $H' \mapsto (H' \wedge F, H' \wedge G)$, et de même pour une
configuration de figures deux à deux disjointes.}
\end{quote}

C'est clair en termes de $\Phi$, dit la page : $\overline{F \cup G}$ est la
réunion disjointe de $\overline{F}$ et $\overline{G}$, entre lesquelles il n'y
a aucune relation d'ordre, et ses parties fermées sont les couples de parties
fermées\footnote{La page écrit $H \cap F$, $H \cap G$ pour $H' \cap F$, $H' \cap
G$.}. D'où la forme définitive (page 57) : la donnée de $(\mathfrak{F}, \leq)$
équivaut à celle d'un ensemble ordonné $\mathcal{M}$ et d'un ensemble $\Phi$ de
parties fermées vers le bas, a) stable par passage aux parties fermées,
b) contenant les $\mathcal{M}_{\leq X}$, c) contenant la partie vide. Et les
parties fermées vers le bas de $\mathcal{M}$ correspondent bijectivement aux
parties de $\mathfrak{F}$ fermées vers le bas et stables par les bornes
supérieures qui existent, par $\mathfrak{F}' \mapsto \mathfrak{F}' \cap
\mathcal{M}$ et $\mathcal{M}' \mapsto \{F \mid \overline{F} \subset
\mathcal{M}'\}$.

\begin{quote}
\textbf{C 3} (page 58). \emph{Si toute strate de $F$ est compatible à toute
strate de $G$, alors $F$ et $G$ sont compatibles.}
\end{quote}

La réciproque est évidente ; une figure est compatible à la réunion d'une
configuration si et seulement si elle l'est à chacun de ses
termes\footnote{Une première version de C 3, sur trois figures deux à deux
compatibles, est barrée en haut de la page 58.}. La structure est alors portée
par $(\mathcal{M}, \leq, R)$, où $R$ est la compatibilité des multistrates,
réflexive, symétrique et héréditaire : $(X, Y) \in R$, $X' \leq X$, $Y' \leq Y$
entraînent $(X', Y') \in R$. Une figure est \emph{finie} si elle a un nombre
fini de strates, \emph{de type fini} si elle est la borne supérieure d'un nombre
fini de strates ; les strates maximales d'une figure finie non vide sont ses
\emph{composantes irréductibles}.

\begin{quote}
\textbf{Proposition} (pages 59 à 61). \emph{La donnée de $(\mathfrak{F}, \leq)$
satisfaisant C 1, C 2, C 3, toutes les figures étant de type fini, équivaut à
celle de $(\mathcal{M}, \leq, R)$, où $R$ est réflexive, symétrique et
héréditaire et où toute partie fermée majorée de $\mathcal{M}$ est de type fini
--- par exemple si les $\mathcal{M}_{\leq X}$ sont finis. Les figures sont
alors les parties de $\mathcal{M}$ fermées, de type fini, et formées
d'éléments deux à deux compatibles.}
\end{quote}

La page 61 en donne la forme finie : les $\mathcal{M}_{\leq X}$ finis, les
figures étant les parties fermées finies dont les éléments sont deux à deux
compatibles\footnote{La page 60 écrit « $A \in \mathcal{M}$ » pour $A \subset
\mathcal{M}$, et la marge renvoie à « Cont 2 de p. 33 », c'est-à-dire à Fig 2.
Le chapitre VI (dossier 156-6, page 16) se référera à ces conditions C 1 à C 3
comme à ce que GF IV avait « explicité en long et en large ».}.

\pagerange{61}{65}
\section*{XII. Composantes connexes (pages 61 à 65)}

Une figure $F$ est \emph{connexe} si $F = F' \amalg F''$, avec $F'$, $F''$
disjointes, entraîne $F' = F$ ou $F'' = F$ ; \emph{$0$-connexe} si de plus elle
n'est pas vide. Une \emph{composante connexe} de $F$ est une sous-figure
$0$-connexe $F_{0}$ telle que $F = F_{0} \amalg F_{1}$. Par la bijection de la
page 57, une sous-figure connexe de $F' \amalg F''$ est contenue dans l'une
des deux, et deux composantes connexes sont égales ou disjointes\footnote{Le
lemme est d'abord énoncé pour « une sous-figure connexe de $G$ », pour $F$. Un
Corollaire 2 sur les familles de composantes, et l'aveu qui le suit --- « Je
n'ai pas pu déduire que [toute] figure soit [réunion de ses] comp.\
connexes » sauf pour les figures finies ---, sont barrés.}.

\begin{quote}
\textbf{Proposition} (pages 62 et 63). \emph{$F_{0} \leq F$ est une composante
connexe de $F$ si et seulement si $\overline{F_{0}}$ est une composante
connexe de l'ensemble ordonné $\overline{F}$, et $F$ est la réunion disjointe
de ses composantes connexes.}
\end{quote}

Une décomposition $F = F' \amalg F''$ est en effet une partition de
$\overline{F}$ en deux parties fermées sans relation d'ordre entre elles,
c'est-à-dire en deux réunions de composantes du graphe de comparabilité de
$\overline{F}$, et réciproquement. La Proposition donne ainsi ce que le
corollaire biffé n'obtenait pas.

Soit $\mathcal{M} = \coprod_{i \in I} \mathcal{M}_{i}$ la décomposition de
l'ensemble ordonné $\mathcal{M}$ en composantes connexes, et, pour une figure
$F$, $F_{i}$ la borne supérieure de ses strates situées dans
$\mathcal{M}_{i}$. Alors $F$ est la somme disjointe des $F_{i}$, presque toutes
vides si $F$ est de type fini ; l'ensemble $\mathfrak{F}_{i}$ des figures
contenues dans $\mathcal{M}_{i}$ satisfait C 1, C 2, et C 3 si $\mathfrak{F}$
le fait ; et $F \mapsto (F_{i})$ est une injection $\mathfrak{F} \to \prod_{i}
\mathfrak{F}_{i}$ qui respecte l'ordre, les bornes supérieures et les bornes
inférieures non vides. Pour que les familles presque nulles soient atteintes, il
faut et il suffit de poser :

\begin{quote}
\textbf{C 4} (page 64). \emph{Deux multistrates situées dans deux composantes
connexes distinctes de $\mathcal{M}$ sont compatibles.}
\end{quote}

L'image est alors formée des familles presque nulles (Corollaire 1), c'est tout
$\prod \mathfrak{F}_{i}$ si $I$ est fini (Corollaire 2), et, si toutes les
figures sont de type fini (resp.\ finies), il en est de même des
$\mathfrak{F}_{i}$ et l'image est exactement l'ensemble des familles presque
nulles (Corollaire 3 et Lemme de la page 65). La page propose aussi un
\textbf{C 4 bis} renforcé, où toute famille $(F_{i}) \in \prod \mathfrak{F}_{i}$
provient d'une figure\footnote{C 4 bis et l'hypothèse que toutes les figures
sont de type fini s'excluent dès que $I$ est infini ; ce sont deux options, non
deux axiomes à cumuler.}. On note $\pi_{0}(F)$, $\pi_{0}(\mathcal{M})$, et
$\pi_{0}(\mathcal{C})$ pour la contrée, en avertissant que ce dernier n'est pas
le $\pi_{0}$ de l'ensemble ordonné $(\mathfrak{F}, \leq)$, qui est connexe
puisqu'il a un plus petit élément, mais celui qui tient compte des bornes
supérieures\footnote{L'exposant de $\pi_{0}(\mathfrak{F})$ est douteux, et le
troisième terme du triplet $(\mathfrak{F}, \leq, \ldots)$ n'est pas lu.}.

\pagerange{65}{68}
\section*{XIII. Supports et cosupports (pages 65 à 68)}

Une figure est disjointe de chaque élément d'une famille $\Phi$ de figures si
et seulement si ses strates le sont (C 3). On pose
\[
\mathrm{cosupp}_{\mathfrak{F}}(\Phi) = \{F \in \mathfrak{F} \mid F \
\text{disjointe de tout}\ H \in \Phi\}, \qquad
\mathrm{cosupp}_{\mathcal{M}}(\Phi) = \mathcal{M} \cap
\mathrm{cosupp}_{\mathfrak{F}}(\Phi),
\]
qui se déterminent l'un l'autre ($F$ est dans le premier si et seulement si
$\overline{F}$ est contenu dans le second), et
$\mathrm{supp}(\Phi) = \mathrm{cosupp}(\mathrm{cosupp}(\Phi))$, le
\emph{support}. Une partie est \emph{saturée} si elle est égale au cosupport de
son cosupport ; toute partie est contenue dans une plus petite partie saturée,
son support ; les parties saturées de $\mathfrak{F}$ et de $\mathcal{M}$ se
correspondent par $\mathcal{M}_{0} = \mathfrak{F}_{0} \cap \mathcal{M}$, et le
cosupport échange les parties saturées en renversant l'inclusion\footnote{En
termes actuels : la relation symétrique « être disjoints » définit une
connexion de Galois de $\mathfrak{P}(\mathcal{M})$ avec lui-même --- une
\emph{polarité} au sens de G.~Birkhoff (\emph{Lattice Theory}, 1940) ---, et
les parties saturées en sont les fermés. Le nom est le nôtre.}.

L'ensemble $\Sigma$ des parties saturées, les \emph{supports virtuels}, a des
bornes inférieures quelconques (les intersections) et des bornes supérieures
quelconques (le saturé de la réunion). « Je présume qu'on a en fait une algèbre
de Boole (\emph{ultra-stonienne}, à cause des Sup quelconques) », dont l'espace
de Stone serait totalement discontinu et ultra-stonien\footnote{Comme aucune
multistrate n'est disjointe d'elle-même, le cosupport est une orthocomplémentation
de $\Sigma$, qui est donc un treillis complet orthocomplémenté ; il n'est pas
distributif en général pour une relation symétrique quelconque. Quand $R$ est
la relation toujours vraie, « disjointes » signifie « sans minorant commun »,
et $\Sigma$ est l'algèbre des ouverts réguliers de l'ensemble ordonné
$(\mathcal{M}, \leq)$ --- celle qu'on attache à un ensemble ordonné en théorie
du forcing ---, qui est bien une algèbre de Boole complète. L'espace de Stone
d'une algèbre de Boole complète est extrêmement discontinu (M.~H.~Stone, 1937 ;
A.~Gleason, 1958) ; c'est vraisemblablement ce que la page appelle
« ultra-stonien », plutôt que les espaces « hyperstoniens » de J.~Dixmier
(1951), qui sont autre chose. La présomption reste non démontrée.}.

Tout cela ne dépend que de $(\mathcal{M}, \leq, R)$, et non de l'ensemble des
figures : le cosupport d'une partie $\mathcal{M}'$ de $\mathcal{M}$ est
\[
\mathrm{cosupp}_{\mathcal{M}}(\mathcal{M}') = \{X \in \mathcal{M} \mid (X, X')
\in R \ \text{et}\ \mathcal{M}_{\leq X} \cap \mathcal{M}_{\leq X'} = \emptyset \
\text{pour tout}\ X' \in \mathcal{M}'\}.
\]
La page écrit seulement $(X, X') \in R$\footnote{Ce qui définirait l'ensemble
des multistrates compatibles à $\mathcal{M}'$, et non son cosupport. La
disjonction s'exprime bien en termes de $(\mathcal{M}, \leq, R)$, ce que la page
veut montrer, mais demande les deux conditions.}.
Exemple : par C 4, une réunion de composantes connexes de $\mathcal{M}$ est
saturée, son cosupport étant la réunion des autres composantes.

\pagerange{69}{75}
\section*{XIV. La subdivision et le raffinement (pages 69 à 75)}

« J'en viens à la relation d'ordre $\preccurlyeq$ (subdivision), de nature très
différente. » Elle n'a ni plus petit ni plus grand élément, deux subdivisions
d'une même figure n'ont pas toujours de borne supérieure, mais parfois une
borne inférieure. Deux subdivisions de $F$ qui admettent une subdivision
commune reçoivent un nom (non pas « compatibles », le mot étant pris), et la
contrée est dite \emph{modérée} quand deux subdivisions d'une même figure en
admettent toujours une commune\footnote{Le mot souligné est surchargé et la
page dit, pour « modérée », que deux subdivisions sont toujours « comparables »,
lu avec doute. La note de marge, avec son schéma $F''' \preccurlyeq F'
\preccurlyeq F$, $F''' \preccurlyeq F'' \preccurlyeq F$, et l'exemple des
complexes simpliciaux, où deux subdivisions d'un même complexe ont toujours une
subdivision commune, font lire la condition de subdivision commune ; deux
subdivisions ne sont presque jamais comparables.}. Les propriétés qu'il a en
vue portent sur les rapports de $\preccurlyeq$ avec $\leq$.

\begin{quote}
\textbf{C 5} (page 69). \emph{Si $F' \preccurlyeq F$, une figure est disjointe de
$F$ si et seulement si elle l'est de $F'$} (il suffit de le demander pour les
multistrates).
\end{quote}

Autrement dit, une subdivision a même support et même cosupport ; en
conséquence, toute partie saturée de $\mathfrak{F}$ est fermée pour
$\preccurlyeq$ aussi, donc une sous-contrée au sens suivant :

\begin{quote}
\textbf{Définition 1} (page 71). Une \emph{sous-contrée} de $\mathcal{C} =
(\mathfrak{F}, \leq, \preccurlyeq)$ est une partie fermée pour $\leq$ et stable
par bornes supérieures --- donc de la forme $\{F \mid \overline{F} \subset
\mathcal{M}'\}$ pour une partie fermée $\mathcal{M}'$ de $\mathcal{M}$ --- et
fermée pour $\preccurlyeq$.

\textbf{C 6} (page 71). \emph{Soient $F' \preccurlyeq F$ et $G \leq F$.} (a)
\emph{Il existe une unique figure $G'$ avec $G' \preccurlyeq G$ et $G' \leq
F'$ ;} (b) \emph{les strates de $G'$ sont les strates de $F'$ qui raffinent
$G$} ; (c) \emph{$G \mapsto G'$ commute aux bornes supérieures.}
\end{quote}
\[
\begin{array}{ccc}
G & \leq & F \\
\succcurlyeq & & \succcurlyeq \\
G' & \leq & F'
\end{array}
\]
$G'$ est la \emph{subdivision de $G$ induite} par $F'$, ou la sous-figure de
$F'$ induite par $G$\footnote{L'énoncé nomme $H$ la figure que son carré appelle
$G'$ ; (c) est dans la marge, reliée à C 6 par une flèche. Le (b) fait appel au
raffinement $\ll$, qui ne sera défini qu'à la page 73 ; il s'entend avec la
définition de la page 73.}.

\begin{quote}
\textbf{C 7}$_{0}$ (page 72). \emph{Si $G \leq F$ et $G' \preccurlyeq G$, il
existe $F' \preccurlyeq F$ qui induit $G'$.}
\end{quote}

L'énoncé symétrique, où l'on échange $\leq$ et $\preccurlyeq$, est faux : les
sous-figures d'une subdivision ne sont pas toutes induites par des sous-figures
de $F$. D'où la dissymétrie que la page relève : « $G$ est sous-figure d'une
subdivision de $F$ » est conséquence de « $G$ est subdivision d'une sous-figure
de $F$ » (par C 7$_{0}$), mais non l'inverse --- un petit segment intérieur à
un triangle est sous-figure d'une subdivision du triangle sans être subdivision
d'aucune de ses sous-figures\footnote{Ce passage de la page 72 est barré, mais
l'énoncé en est juste ; l'exemple est le nôtre. Un C 7 qui choisirait $F'$
canoniquement est annoncé et non écrit.}.

\begin{quote}
\textbf{Définition} (page 73). $G$ \emph{raffine} $F$, et l'on écrit $G \ll F$,
s'il existe $F'$ avec $G \leq F' \preccurlyeq F$.
\end{quote}

La relation $\ll$ est réflexive et impliquée par $\leq$ et par $\preccurlyeq$.
Elle est transitive : si $F \leq G' \preccurlyeq G$ et $G \leq H'
\preccurlyeq H$, C 7$_{0}$ appliqué à $G \leq H'$ et $G' \preccurlyeq G$ donne
$H'' \preccurlyeq H'$ avec $G' \leq H''$, d'où $F \leq H'' \preccurlyeq
H$\footnote{La page dit que la démonstration « utilise C 6, ni C 5 ni C 7 », et
son schéma, tel qu'il est lu, passe par une décomposition $G \preccurlyeq
\overline{H} \leq H$ que le passage barré de la page 72 déclarait non
toujours disponible. La démonstration par C 7$_{0}$ est la nôtre ; c'est
d'ailleurs celle que la page 83 emploie.}. Pour l'antisymétrie, la page se
ramène au

\begin{quote}
\textbf{Lemme} (page 73). \emph{Si $F \preccurlyeq K \leq F$, alors $K = F$.}
\end{quote}

et le déduit de C 5 et d'un corollaire : si $K < F$, il existe une multistrate
disjointe de $K$ et non de $F$, de sorte que $K$ et $F$ n'ont pas même support,
alors qu'une subdivision a même support. Mais ce corollaire n'est pas
« formel » en termes de C 1 à C 4 (page 74) : si $\mathfrak{F}$ est l'ensemble
de toutes les parties fermées d'un ensemble ordonné $\mathcal{M}$ qui a un plus
petit élément, deux figures non vides ne sont jamais disjointes, et toutes ont
le même support. Avec $\preccurlyeq$ discrète, C 6 et C 7 sont alors
tautologiques (page 75). « On voit donc qu'il manque un axiome »\footnote{La page
75 renvoie à Cont 4 (page 35) et à Cont 7 (page 41) comme à la forme cherchée ;
un premier essai sur le raffinement induit y est barré.}.

\pagerange{75}{80}
\section*{XV. Raffinement induit et plus petite strate (pages 75 à 80)}

« Il faut développer la notion de raffinement induit. »

\begin{quote}
\textbf{Proposition} (page 76). \emph{Soient $G \ll F$ et $F' \leq F$. Il
existe une plus grande sous-figure $G'$ de $G$ qui raffine $F'$, le
\textbf{raffinement de $F'$ induit} par $G$ ; ses strates sont les strates de
$G$ qui raffinent $F'$, et $G'$ est la borne inférieure de $G$ et de $F'$ pour
$\ll$.}
\end{quote}
\[
\begin{array}{ccc}
G & \ll & F \\
\geq & & \geq \\
G' & \ll & F'
\end{array}
\]

La marge ajoute qu'on posera $G' = \emptyset$ si et seulement si $G$ est
disjointe de $F'$, seul le sens direct n'étant pas évident\footnote{La page
écrit « une plus grande sous-figure $F'$ de $G$ » pour $G'$. Une première
démonstration, barrée, contient deux autres confusions de lettres.}. Écrivant
$G \leq \overline{F} \preccurlyeq F$ et appliquant C 6 à $\overline{F}
\preccurlyeq F$ et $F' \leq F$, on obtient $\overline{F}' \preccurlyeq F'$,
$\overline{F}' \leq \overline{F}$, et $G' = G \wedge \overline{F}'$, qui
raffine $F'$. Pour la propriété de borne inférieure, il faut montrer que $K \ll
G$ et $K \ll F'$ entraînent $K \ll G'$, et l'on se ramène aux deux cas $G
\preccurlyeq F$ et $G \leq F$. Le \textbf{Lemme 1} (page 77) traite le premier
par C 6 (b)\footnote{En utilisant, sans le dire, la transitivité de
l'induction que la page 49 notait au quatrième départ.} ; le \textbf{Lemme 2}
(page 78), le second, commence par $K \leq \overline{G} \preccurlyeq G$ et
s'arrête : « J'abandonne ici ! Il [vaut] mieux se reporter page 51 »\footnote{La
lecture de « J'abandonne ici » est douteuse. La page 51 est le Corollaire du
quatrième départ, qui établissait la même propriété dans l'autre cadre. Une
note de marge affirme encore que $F' \mapsto G'$ commute aux bornes supérieures
et aux bornes inférieures quelconques.}. La propriété d'Inf reste ainsi non
démontrée dans ce départ.

\begin{quote}
\textbf{Proposition} (pages 78 et 79). \emph{Soient $F$ une figure et $X \in
\mathcal{M}$ avec $X \ll F$.} a) \emph{Il existe une strate $Y$ de $F$ avec
$X \ll Y$.} b) \emph{Parmi ces strates, il y en a une plus petite : celle dont
aucune sous-multistrate stricte n'est raffinée par $X$.}
\end{quote}

a) Comme $F = \mathrm{Sup}_{i} Y_{i}$ sur ses strates, C 6 (c) donne $X =
\mathrm{Sup}_{i} G_{i}$, où $G_{i}$ est la sous-figure de $X$ induite par
$Y_{i}$ ; $X$ étant irréductible, $X = G_{i}$ pour un $i$, et $X \ll Y_{i}$.
b) Un élément minimal existe si les multistrates sont de dimension
combinatoire finie, « on y reviendra » ; il est le plus petit, parce que si
$X \ll Y$ et $X \ll Y'$, la sous-figure de $X$ induite par $Y \wedge Y'$ est
encore $X$, de sorte que $X$ raffine une strate de $Y \wedge Y'$, majorée par
$Y$ et donc égale à $Y$ par minimalité\footnote{La page caractérise $Y$ par
« la sous-figure de $X$ induite par $Y$ est $\neq X$ » ; la page 79 dit le
contraire, et c'est elle qu'on suit : la sous-figure induite par $Y$ est $X$,
celle qu'induit le bord de $Y$ ne l'est pas. L'argument d'unicité invoque « les
propriétés fonctorielles des Inf pour $\ll$ », c'est-à-dire la commutation de
l'induction aux bornes inférieures, affirmée en marge de la page 78 et non
démontrée. La page conclut $Y = Y'$ pour un $Y'$ minimal ; pour un $Y'$
quelconque on obtient $Y \leq Y'$.}.

\begin{quote}
\textbf{Corollaire} (page 80). \emph{Si $F \ll G$, l'application $\varphi :
\overline{F} \to \overline{G}$ qui associe à une strate $X$ de $F$ la plus
petite strate de $G$ qu'elle raffine est croissante.}
\end{quote}

C'est l'application que la page 42 cherchait sans pouvoir l'obtenir, et
l'analogue de celle que le chapitre V construit le même jour pour les familles
admissibles\footnote{Dossier 156-5, pages 19 et 20, daté du 14 juin.}.
« Jusqu'à maintenant », remarque la page, les axiomes C 5, C 6, C 7 sont
tautologiques quand $\preccurlyeq$ est l'égalité : « Il serait temps d'en venir
à des propriétés plus spécifiques ! »

\pagerange{80}{85}
\section*{XVI. 14 juin : l'antisymétrie du raffinement (pages 80 à 85)}

Le 14 juin, il renomme $\mathcal{M}_{\ll F}$ l'ensemble des multistrates qui
raffinent $F$, « pour éviter confusion avec notion d'adhérence »\footnote{Il
note cet ensemble $\widetilde{F}$, au lieu de $\overline{F}$ dans l'« ancien
C 4 » de la page 47.}, et énonce un corollaire de la Proposition de la page 78
qui reprend, dans ce départ, la relation encadrée $G' \cap F \subset G$ de la
page 47 :

\begin{quote}
\textbf{Corollaire} (pages 80 à 82). \emph{Si $G \leq F$, toute strate de $F$
qui raffine $G$ est une strate de $G$ :} $\overline{F} \cap \mathcal{M}_{\ll G} =
\overline{G}$.
\end{quote}

La formule de la page est « soient $G' \preccurlyeq G \preccurlyeq F$, alors
$\widetilde{G}' \cap \widetilde{F} \subset \widetilde{G}$ » ; lue ainsi, elle
est triviale, puisque $\ll$ est transitive, et les deux signes sont repassés à
l'encre. On suit la paraphrase qui la suit et la démonstration, qui partent
d'une strate $X$ de $F$\footnote{La reconstitution est la nôtre. Pour deux
figures $X \leq F$ et $G \leq F$, elles sont compatibles, et le corollaire est
alors la forme, dans ce départ, du Corollaire 3 de la page 48 : une
multistrate compatible à $G$ qui la raffine en est une strate.}. La
démonstration (pages 80 à 82) passe par les sous-figures induites sur $X$ par
$G$ et par une subdivision de $F$, se ramène à prouver $X \leq G_{X}$, et
achève par l'axiome que voici, marqué « provisoire ! » :

\begin{quote}
\textbf{C 8 (0)} (page 82). \emph{Pour $X \in \mathcal{M}$ et $\Phi, \Psi \in
\mathfrak{F}$ :} (a) \emph{$X \leq \Phi \preccurlyeq X$ entraîne $X = \Phi$ ;}
(b) \emph{$X \preccurlyeq \Psi \leq X$ entraîne $X = \Psi$.}
\end{quote}

Il est nécessaire si $\ll$ doit être un ordre : dans les deux cas, chacun des
deux termes raffine l'autre\footnote{Le haut de la page 81 et un premier énoncé
de C 8, au bas, sont barrés ; plusieurs relations verticales des schémas de la
page 81 sont de lecture douteuse. La page 82 ajoute que C 8 est aussi
« suffisant », mot lu avec doute.}.

\begin{quote}
\textbf{Corollaire} (pages 82 et 83). \emph{Si $F \ll G$ et $G \ll F$, alors
$F = G$.}
\end{quote}

On a $F \leq G' \preccurlyeq G \leq F' \preccurlyeq F$. C 7$_{0}$, appliqué à
$G \leq F'$ et $G' \preccurlyeq G$, donne $H \preccurlyeq F'$ avec $G' \leq H$,
d'où
\[
(*) \qquad F \leq H \preccurlyeq F .
\]
De $F = H$ viendraient $G' = F$ et $F' = F$, puis
\[
(**) \qquad F \preccurlyeq G \leq F ,
\]
et de là $F \leq G$ par le corollaire de la page 80 appliqué aux strates de $F$
--- chacune raffine $G$, puisque $F \preccurlyeq G$ ---, donc $F = G$. Mais
l'égalité $F = H$ dans $(*)$ est tirée de C 8 (a), qui ne vaut que pour une
multistrate ; pour une figure quelconque elle n'est pas
établie\footnote{La note de marge de la page 83, lue par fragments, semble le
relever : « … de $F$ et $H$ … $H \preccurlyeq F \leq H$ ! ». « On a essayé ici
de se débrouiller aussi bien que possible, sans axiomes et raisonnements de
disjonction, genre C 5. »}.

\subsection*{La structure dégagée}

La page 84 en tire une structure abstraite.

\textbf{Remarque} (pages 84 et 85). \emph{Soit $\mathfrak{F}$ un ensemble muni de
deux ordres $\leq$ et $\preccurlyeq$ tels que :}
\begin{enumerate}
\item[a)] \emph{si $G \leq F$ et $F' \preccurlyeq F$, il existe un unique $G'$
avec $G' \preccurlyeq G$ et $G' \leq F'$ ;}
\item[b)] \emph{si $G \leq F$ et $G' \preccurlyeq G$, il existe un $F'$, non
nécessairement unique, avec $F' \preccurlyeq F$ et $G' \leq F'$ ;}
\item[c)] \emph{$F \leq G \preccurlyeq F$ entraîne $F = G$.}
\end{enumerate}
\emph{Alors (i) la relation $G \ll F$, « il existe $F'$ avec $G \leq F'
\preccurlyeq F$ », est un ordre ; (ii) le carré de a) est cartésien pour $\ll$,
c'est-à-dire $G' = \mathrm{Inf}_{\ll}(F', G)$ ; (iii) un carré cartésien pour
$\geq$ l'est pour $\ll$.}

a) est C 6, b) est C 7$_{0}$, c) est la forme de C 8 pour toutes les figures ;
la condition c) s'écrit aussi $G \preccurlyeq F \leq G \Rightarrow F = G$, qui
est la même aux noms près. (i) est démontré par ce qui précède : la
transitivité par b), l'antisymétrie par b) et c), comme aux pages 82 et 83.
(ii) et (iii) sont affirmés sans démonstration --- c'est la propriété d'Inf
abandonnée à la page 78 ---, et une note de marge doute encore de la bonne
définition\footnote{Les relations de (iii) sont surchargées et de lecture
douteuse ; la note de marge n'est lue que par fragments.}.

Il faudrait, conclut la page 85, séparer les axiomes des contrées en paquets :
un paquet sur $\leq$ seul (C 1 à C 4) ; un paquet de propriétés « formelles »
de $\leq$ et $\preccurlyeq$, sans questions de support ni de disjonction, qui
assurent que les bornes supérieures commutent à l'induction sur une
subdivision ; des axiomes de recollement des subdivisions ; un paquet d'axiomes
de disjonction.

\pagerange{86}{88}
\section*{XVII. Recollement, lieux, ombres (pages 86 à 88)}

\begin{quote}
\textbf{C 8 --- Axiome de recollement des subdivisions} (page 86). \emph{Soit $F
= \mathrm{Sup}_{i} F_{i}$, et pour tout $i$ une subdivision $F'_{i}$ de $F_{i}$,
telles que $F'_{i}|F_{i} \wedge F_{j} = F'_{j}|F_{i} \wedge F_{j}$. Alors
$\mathrm{Sup}_{i} F'_{i}$ existe et subdivise $F$.}
\end{quote}

C'est le C 7 de la page 49 transporté dans ce départ, qui « complète C 6 (c) »,
et le nom « C 8 » est réemployé pour lui\footnote{Le C 8 (0) de la page 82, le
C 8 barré de la page 78, et le Cont 8 de la page 43 sont trois autres énoncés ;
d'où la règle, suivie ici, de toujours citer un axiome avec sa page.}. Pour une
somme disjointe $F = \coprod F_{i}$, il donne
\[
\mathrm{Subdiv}(F) \simeq \prod_{i} \mathrm{Subdiv}(F_{i}).
\]

Puis le titre « Axiomes de disjonction et de support », et : « pour formuler
les axiomes sous forme aussi forte et sensible que possible, c'est le moment
d'introduire les lieux »\footnote{Une tentative, barrée au bas de la page 86 et
en haut de la page 87, montrait l'équivalence, pour une figure non vide, de la
minimalité pour $\preccurlyeq$ et de la minimalité pour $\ll$ parmi les figures
non vides.}.

\begin{quote}
\textbf{Définition} (page 87). Un \emph{lieu} de la contrée est une figure non
vide minimale pour $\ll$ parmi les figures non vides. On note $\mathcal{L}$
l'ensemble des lieux.
\end{quote}

Un lieu $x$ est une multistrate --- ses strates sont des figures non vides qui
le raffinent, donc lui sont égales --- et c'est un élément minimal de
$\mathcal{M}$ pour $\leq$. La page note que la réciproque n'est pas vraie avec
les axiomes posés jusqu'ici, et le sera pour les « contrées sans
bords »\footnote{La notion n'est pas définie dans le dossier. Dans le modèle
topologique des pages 1 à 19, la réciproque est fausse même pour un espace sans
bord : un cercle, une seule strate, est minimal pour $\leq$, mais chacun de ses
points le raffine (page 19). « Sans bords » doit donc vouloir dire autre chose
que l'absence de bord des variétés.}.

L'\emph{ombre} d'une figure $F$ est $\mathcal{L}_{F} = \{x \in \mathcal{L} \mid
x \ll F\}$, et sa \emph{multiombre} est la famille de parties de $\mathcal{L}$
\[
\widehat{\mathcal{L}}_{F} = \{\mathcal{L}_{X} \mid X \in \overline{F}\} \subset
\mathfrak{P}(\mathcal{L}).
\]
Ce sont les objets du premier et du troisième départ, et l'on garde leurs
notations\footnote{La page note l'ombre $|F|$ et la multiombre
$\widehat{F}$.}.
On a $\mathcal{L}_{F} = \bigcup_{X \in \overline{F}} \mathcal{L}_{X}$ : c'est
la Proposition de la page 78, qui fournit même une plus petite strate de $F$
raffinée par un lieu donné.

\begin{quote}
\textbf{Proposition} (page 88). \emph{La multiombre d'une figure est une famille
admissible de parties de $\mathcal{L}$ :} a) \emph{pour $A \in
\widehat{\mathcal{L}}_{F}$, posant $\partial A = \bigcup_{B \subsetneq A} B$, on
a $A^{\circ} = A \smallsetminus \partial A \neq \emptyset$ ;} b) \emph{pour $A,
B \in \widehat{\mathcal{L}}_{F}$, $A \cap B$ est la réunion des $C \in
\widehat{\mathcal{L}}_{F}$ contenus dans $A \cap B$.}
\end{quote}

C'est, dans le cinquième départ, le Scholie des pages 23 et 36 ; mais ici la
démonstration ne suit pas, et une note de marge prévient qu'elle utilise
quelque chose qui n'est pas lu. Le dossier s'arrête sur cet énoncé\footnote{La
deuxième mouture de l'Analysis situs (chapitre VI, dossier 156-6) en fera un
axiome, « At 8 » de sa page 45, et y reprendra la définition des lieux comme
éléments minimaux pour $\ll$ (sa page 54) ; c'est ce qu'en disent les notes de
sa transcription.}. La boucle est bouclée sans être fermée : la page 24
retrouvait la contrée à partir des ombres de ses multistrates ; la page 88
voudrait montrer que ces ombres ont encore la forme d'une figure, cette fois à
partir d'un cadre où les lieux ne sont plus primitifs, mais les derniers venus.

\end{document}
